\documentclass[aps,pra,reprint,superscriptaddress,longbibliography]{revtex4-2}
\usepackage[T1]{fontenc}
% Conditional typography fallbacks for incomplete offline font caches.
\IfFileExists{ec-lmr5.tfm}{}{%
 \IfFileExists{t1lmr.fd}{%
  \input{t1lmr.fd}%
  \DeclareFontShape{T1}{lmr}{m}{n}
   {<-6.5>ec-lmr6 <6.5-7.5>ec-lmr7 <7.5-8.5>ec-lmr8
    <8.5-9.5>ec-lmr9 <9.5-11>ec-lmr10 <11-15>ec-lmr12
    <15->ec-lmr17}{}%
 }{}%
}
\usepackage{amsmath,amssymb,bm,booktabs,graphicx}
\IfFileExists{eufm7.pfb}{}{%
 \input{ueuf.fd}%
 \DeclareFontShape{U}{euf}{m}{n}{<->eufm10}{}%
}
\IfFileExists{cmbx7.pfb}{}{%
 \input{ot1cmr.fd}%
 \DeclareFontShape{OT1}{cmr}{bx}{n}
  {<5>cmbx5 <6>cmbx6 <7>cmbx8 <8>cmbx8 <9>cmbx9
   <10>cmbx10 <10.95>cmbx10 <12>cmbx12 <14.4>cmbx12
   <17.28>cmbx12 <20.74>cmbx12 <24.88>cmbx12}{}%
}
\usepackage[colorlinks=true,linkcolor=blue,citecolor=blue,urlcolor=blue]{hyperref}
\newcommand{\Tr}{\operatorname{Tr}}
\newcommand{\dd}{\mathrm d}
\newcommand{\ee}{\mathrm e}
\newcommand{\id}{\mathbb I}
\newcommand{\supp}{\operatorname{supp}}
\newcommand{\norm}[1]{\left\lVert#1\right\rVert}
\newcommand{\ket}[1]{\lvert#1\rangle}
\newcommand{\bra}[1]{\langle#1\rvert}
\newcommand{\ip}[2]{\langle#1,#2\rangle}
\newcommand{\one}{\mathbf 1}
\begin{document}
\title{Transverse information beyond thermal scalar limits in colliding layers}
\author{R. Arun Chandru}
\email{arun.chandru@gmail.com}
\date{\today}
\begin{abstract}
Thermal scalar observables can approach the same limits while coherent
layer measurements retain information about a collapsing geometry.
We establish this separation for the actual Green operators of finitely
many weighted curves approaching a common curve. A joint separation--temperature
limit determines the partition function, normalized heat state, and
transverse population, including their high-frequency tails. In a subsequent
high-temperature limit, the one-excitation entropy differences vanish,
whereas $M$ equally weighted layers remain distinguishable from their
collapsed reference with trace distance tending to $1-1/M$. The finite
crossover has a strict counting deficit, and its low-temperature
transverse population detects rotation of the lowest spectral branch.
We then distinguish controlled one-excitation preparation from a full
number-conserving Bose ensemble. Calibrated occupation thresholds tolerate
fixed sorting errors; a common infrared filter yields exponentially small
error with an explicit stopband-leakage budget. Exact circular Green blocks
give a finite instance with average error below $0.0032$ under
specified collection loss, sorting errors, and dark counts. Preparation,
occupied modes, and the common-vacuum assumption are also included.
\end{abstract}
\maketitle

\section{Introduction}

Thermal measurements compress a spectrum. Partition functions, mean
energies, and entropies summarize the occupations of its modes, but a
measurement that mixes those modes also probes their orientation in the
accessible Hilbert space. The distinction is particularly consequential
when a geometric limit and a thermal limit occur together. Modes that
disappear at every fixed spectral index can still carry a finite fraction
of the information in a temperature-dependent spectral window. The
question addressed here is whether this mechanism survives for an actual
continuum geometry, with a controlled measurement and preparation cost.

Consider mass supported on several close curves in a planar carrier.
Their common displacement and relative displacements define longitudinal
and transverse layer combinations. At fixed energy the curves merge,
but wavelengths comparable with their separation resolve the layers.
We compute the resulting joint limit rather than first replacing the
layers by a single curve. The limiting object is a positive matrix-valued
symbol whose off-diagonal entries are decaying exponentials of the
normal separations. Its eigenvalues determine thermal scalar quantities;
its eigenvectors determine the coherent transverse response.

There are close precedents on both sides of this problem. Singular
interactions supported by curves, parallel lines, and colliding supports
have a substantial spectral theory
\cite{KondejKrejcirik2013,KondejKrejcirik2016}. The sharp
one-dimensional Lieb--Thirring inequality for measures
\cite{HundertmarkLiebThomas1998} controls a finite-profile counting
deficit that appears below. Neither singular interactions nor their
exponential Birman--Schwinger kernels are new ingredients. Concentrating
physical mass has also been studied through boundary layers and thin
heavy inclusions \cite{DallaRivaProvenzano2018,Golovaty2022}; their
bulk-mass and boundary regimes differ from the interior pushforward
probability masses used here.
High-index spectral limits under concentrated mass have also been studied
for branched surfaces with Kirchhoff junctions
\cite{GolovatyGomezPerezMartinez2025}. Our limit instead concerns
fixed-total-mass curves that collide and their trace-class thermal
information laws. Matrix semiclassical trace formulas and
nonsmooth matrix-symbol tests provide established analytic context
\cite{AssalDimassiFujii2019,Kuster2017,BollmannMuller2024}.
The singular inverse symbol and both frequency tails still require
the actual-operator comparison proved below. On the
measurement side, optical source discrimination already uses a reference
mode and its orthogonal complement to detect information missed by
direct imaging \cite{Lu2018}. The failure of an ideal dark-port rule
under imperfect mode sorting, and its replacement by a calibrated
threshold, are also established \cite{Schlichtholz2024}.
Background-robust discrimination using reduced spatial-mode
demultiplexing has recently been studied in a weak-source Poisson
regime \cite{Liu2026}; the count bounds below retain full Bose
occupation statistics. Bosonic
determinant identities and passive optical synthesis provide the relevant
counting and implementation tools \cite{Klich2002,Clements2016}.

The contribution is the connection of these ingredients through a
specific continuum limit and its operational consequences. We prove an
all-mode joint law for the actual curved Green operators, including a
trace-norm comparison on a common space. This gives a separation between
thermal scalar agreement and transverse discrimination, a strict
finite-crossover deficit, and a profile-dependent distinction between
excited-branch occupation and rotation of the lowest branch. We then
derive preparation and measurement statements in two explicitly
different ensembles. In the full Bose model a common low-frequency
rejection filter preserves a macroscopic signal while controlling the
largest modal fluctuations; its error exponent has an explicit dependence
on stopband leakage. Finally, a finite circular example replaces an
unspecified asymptotic onset by exact Green matrices and validated
probability bounds.

The physical model is specified at two levels. The continuum operator
describes a static Dirichlet energy with mass on curves. Its quantum
access model uses common, number-conserving modes, which can be programmed
in a finite optical encoding. An ordinary mechanical quantization can
require an additional Bogoliubov transformation because its vacuum
changes with the frequencies. Separating these levels allows the
statistical conclusions to be stated without attributing an unbuilt
measurement architecture to a membrane. The distinction between thermal
scalars and coherence also has antecedents in quantum thermodynamics
\cite{Lostaglio2015}. Here each state is Gibbs-diagonal in its own
Hamiltonian: the comparison concerns different geometric Hamiltonians
in a common layer alignment, not coherence relative to one fixed
Hamiltonian or an additional work-extraction law. Full derivations are included
in the appendices.

\section{Geometry, alignment, and the resolved thermal scale}
\label{sgt:main:model}

Let $\Gamma$ be a smooth embedded closed curve of length $L$ in a bounded
planar Dirichlet carrier $\Omega$. Its fixed tubular neighborhood lies
strictly inside $\Omega$. Write $s$ for arclength, $\gamma_0(s)$ for its
position, and $n(s)$ for a chosen unit normal. Fix distinct heights
$a_1<\cdots<a_M$, positive weights $\alpha_i$ with
$\sum_i\alpha_i=1$, and a smooth positive density $w$ satisfying
$\int_\Gamma w\,\dd s=1$. For sufficiently small $h>0$, the mass measure is
\begin{equation}
 \nu_h=\sum_{i=1}^M\alpha_i
       (s\mapsto\gamma_0(s)+ha_i n(s))_*(w(s)\,\dd s).
 \label{sgt:main:measure}
\end{equation}
Thus the mass is prescribed by pushforward, not by inserting the
arclength Jacobian of each offset curve. Fixed absorbing sets outside
the tube may be incorporated into the carrier.

Restrict the Dirichlet Green kernel to these curves and conjugate by
$\sqrt{\alpha_iw}$. The resulting positive compact inverse $R_h$ acts on
the common space
\begin{equation}
\begin{gathered}
 \mathcal H=L^2(\Gamma,\dd s)\otimes\mathbb C^M,
 \\
 v_i=\sqrt{\alpha_i},\\
 P=vv^*,\\
 Q=I-P.
 \label{sgt:main:alignment}
\end{gathered}
\end{equation}
Equivalently, $L_h=R_h^{-1}$ is the energy operator obtained by minimizing
the bulk Dirichlet energy at fixed curve values, with the stated mass
normalization. The collapsed inverse is $R_0=R_{\mathrm b}\otimes P$.
Its $Q$ complement has zero heat weight: it represents excluded modes,
not additional zero-energy excitations.

Fix a dispersion exponent $\gamma>0$ and set
\begin{equation}
\begin{gathered}
 H_h=L_h^\gamma,\\
 \beta=\tau h^\gamma,\\
 A_h=\ee^{-\beta H_h},\\
 \rho_h=A_h/Z_h,\\
 Z_h=\Tr A_h .
 \label{sgt:main:state}
\end{gathered}
\end{equation}
Here and below the displayed energy and temperature scales are
dimensionless. The mechanical frequency convention corresponds to
$\gamma=1/2$ when $L_h$ is a squared-frequency operator. A reference
state $\rho_{0,h}$ uses $R_0$ at the same $\beta$. Every geometric
parameter, $M$, the positive weights, $\gamma$, and $\tau$ is fixed
before $h\downarrow0$. Limits in $\tau$ are subsequent limits; none of
the conclusions presupposes uniformity when the number of layers grows
or their weights degenerate with $h$.

For $q=|\xi|>0$ define
\begin{align}
 T(q)_{ij}&=\sqrt{\alpha_i\alpha_j}\,
               \ee^{-q|a_i-a_j|},\nonumber\\
 B(s,\xi)&=\frac{w(s)}{2q}T(q),
 &B_0(s,\xi)&=\frac{w(s)}{2q}P,\nonumber\\
 F(s,\xi)&=\ee^{-\tau B(s,\xi)^{-\gamma}},
 &F_0(s,\xi)&=\ee^{-\tau(2q/w(s))^\gamma}P.
 \label{sgt:main:symbol}
\end{align}
The value of the heat function at a zero inverse eigenvalue is zero.
The exponential matrix $T(q)$ is strictly positive for distinct heights.
It contains the entire finite-profile crossover; the remote boundary
contributes lower-order terms on this scale.

\section{What thermal scalar measurements miss}
\label{sgt:main:separation}

Define the phase-space coefficients
\begin{align}
 C(\tau)&=\frac1{2\pi}\int_\Gamma\int_{\mathbb R}
                        \operatorname{tr} F\,\dd\xi\,\dd s,\nonumber\\
 C_Q(\tau)&=\frac1{2\pi}\int_\Gamma\int_{\mathbb R}
                        \operatorname{tr}(QF)\,\dd\xi\,\dd s,\nonumber\\
 C_0(\tau)&=\frac{\Gamma(1+1/\gamma)}{2\pi\tau^{1/\gamma}}.
 \label{sgt:main:C}
\end{align}
The actual operators obey
\begin{gather}
 hZ_h\longrightarrow C(\tau),\\
 hZ_{0,h}\longrightarrow C_0(\tau),\\
 \Tr(Q\rho_h)\longrightarrow p_Q(\tau)=C_Q(\tau)/C(\tau),
 \label{sgt:main:joint}\\
 \frac12\norm{\rho_h-\rho_{0,h}}_1\longrightarrow
\\
 D(\tau)=\frac1{4\pi}\int_\Gamma\int_{\mathbb R}
 \operatorname{tr}\left|\frac F{C(\tau)}-\frac{F_0}{C_0(\tau)}\right|
 \,\dd\xi\,\dd s .
 \label{sgt:main:distance}
\end{gather}
For a nontrivial profile, $0<p_Q(\tau)\le D(\tau)<1$ at every fixed
$\tau>0$. These statements include the low- and high-frequency tails
of the actual Green operators, not only a formal symbol calculation.
Appendix~\ref{sgt:geo:appendix} gives the positive Fourier comparison
and trace argument.

The clearest separation occurs upon subsequently taking $\tau\downarrow0$.
Then $C/C_0\to1$, while
\begin{equation}
 p_Q(\tau)\longrightarrow1-\sum_i\alpha_i^2.
 \label{sgt:main:Qhot}
\end{equation}
For equal masses, the complete trace-distance limit agrees with this
simple layer measurement. The usual binary discrimination relation
\cite{Watrous2018} then gives its operational meaning:
\begin{equation}
 \lim_{\tau\downarrow0}\lim_{h\downarrow0}
 \frac12\norm{\rho_h-\rho_{0,h}}_1=1-\frac1M.
 \label{sgt:main:equal}
\end{equation}
With equal prior probabilities, a single controlled excitation can
therefore be classified with limiting error $1/(2M)$. On $N$ independent
controlled excitations, declaring the split geometry when at least one
$Q$ event occurs gives limiting error $1/(2M^N)$. The latter is the
performance of this specified receiver, without a claim of optimality
over all collective $N$-copy measurements.

This surviving information is compatible with agreement of a substantial
set of thermal scalars. For $p>0$ put
\begin{equation}
 C_p(\tau)=\frac1{2\pi}\int
   \operatorname{tr}\bigl[(\tau B^{-\gamma})^p
                        \ee^{-\tau B^{-\gamma}}\bigr]\,\dd s\,\dd\xi.
 \label{sgt:main:moments}
\end{equation}
At fixed $\tau$, the one-excitation entropy and energy satisfy
\begin{align}
 S(\rho_h)&=\log(1/h)+\log C+C_1/C+o(1),\nonumber\\
 \beta\langle H_h\rangle&\longrightarrow C_1/C,\nonumber\\
 \beta^2\operatorname{Var}_{\rho_h}(H_h)&\longrightarrow
                              C_2/C-(C_1/C)^2.
 \label{sgt:main:scalars}
\end{align}
Every fixed positive-order R\'enyi entropy has an analogous formula.
As $\tau\downarrow0$, the additive entropy differences from the collapsed
reference vanish, and
$C_p/C\to\Gamma(p+1/\gamma)/\Gamma(1/\gamma)$.
Both states also have the same limiting expectations against fixed
smooth spatial multipliers. Thermal scalar agreement here means these
specified limits. It does not mean equality of the spectra or of
arbitrary measurements at finite separation and temperature.

\section{Counting deficit and rotation of the lowest branch}
\label{sgt:main:structure}

The same matrix symbol connects the crossover with a classical
one-dimensional spectral problem. Let $-\kappa_j(g)^2$ be the negative
eigenvalues of
\begin{equation}
 -\frac{\dd^2}{\dd a^2}-g\sum_i\alpha_i\delta_{a_i}
 \quad\hbox{on }\mathbb R,
 \label{sgt:main:delta}
\end{equation}
with absent branches contributing zero. The actual growing-index count is
\begin{equation}
\begin{gathered}
 h\,\#\{\lambda_k(L_h)\le E/h\}\longrightarrow
\\
 \mathcal N(E)=\frac1\pi\int_\Gamma
                \sum_j\kappa_j(w(s)E)\,\dd s
 <\frac{E}{2\pi}
 \label{sgt:main:count}
\end{gathered}
\end{equation}
for $E>0$ and at least two distinct positive-mass layers. The strict
inequality is the finite-profile consequence of the sharp
Lieb--Thirring bound for measures
\cite{HundertmarkLiebThomas1998}; we include its finite-atomic argument
to make the connection transparent. Integrating the counting law gives
$C(\tau)<C_0(\tau)$ at every finite $\tau$. Thus the high-temperature
agreement of leading coefficients coexists with a strict crossover
deficit, not with an exactly invisible spectral change.

There is also a useful distinction within the transverse response.
Define
\begin{gather}
 d_i=\sum_j\alpha_j|a_i-a_j|,\\
 G=\sum_i\alpha_i d_i,\\
 V=\sum_i\alpha_i(d_i-G)^2,\nonumber\\
 I_k=\int_\Gamma w(s)^k\,\dd s,\\
 G_\gamma=\Gamma(1+1/\gamma).
 \label{sgt:main:profile}
\end{gather}
After the joint limit, the large-$\tau$ expansions are
\begin{align}
 C(\tau)&=\frac{G_\gamma}{2\pi}\tau^{-1/\gamma}
 -\frac{G I_2\Gamma(1+2/\gamma)}{4\pi}\tau^{-2/\gamma}
 +O(\tau^{-3/\gamma}),\nonumber\\
 p_Q(\tau)&=\frac{V I_3\Gamma(1+3/\gamma)}{12G_\gamma}
                   \tau^{-2/\gamma}+O(\tau^{-3/\gamma}).
 \label{sgt:main:rotation}
\end{align}
Excited branches of the limiting symbol are exponentially suppressed in
$\tau$, but the lowest eigenvector rotates away from the fixed common
layer vector $v$. The mean-distance statistic $G$ in the first line
is the classical weak-coupling coefficient \cite{Simon1976,Simon2022}.
The second line measures the rotation. Its coefficient
vanishes only for a single layer or two equally weighted layers; in the
latter case the lowest eigenvector remains exactly symmetric and the
transverse response decays exponentially. Replacing $Q$ by the
instantaneous excited-branch projector would therefore remove a genuine
observable effect. The tridiagonal inverse of $T(q)$ gives all finite
branch thresholds and proves these assertions in
Appendix~\ref{sgt:geo:structure}.

\section{Preparing the one-excitation state}
\label{sgt:main:preparation}

The state $\rho_h$ is a normalized one-excitation Gibbs state. Conditioning
a full equilibrium sample on exactly one excitation need not be an
efficient way to prepare it. A uniform spectral estimate for the actual
curves resolves the finite-dimensional requirement:
\begin{equation}
 r_j(R_h)\le K/j,\qquad \lambda_j(L_h)\ge cj.
 \label{sgt:main:weakS1}
\end{equation}
The constants are uniform through the collision at fixed geometric
data. Self-layer regularity and positivity of the full block operator
prove this without assuming uniformly smooth cross-layer kernels.
For fixed trace accuracy $\delta$, the required energy truncation has
\begin{equation}
 J=O\bigl(h^{-1}[1+\log(1/\delta)]^{1/\gamma}\bigr).
 \label{sgt:main:J}
\end{equation}
A common tensor space containing the retained vectors of both hypotheses
has dimension at most $M(M+1)J$ and preserves $P$ and $Q$ exactly.
Selected energy and entropy observables are controlled by their own
weighted tails, rather than inferred from trace distance alone.

On this known finite encoding, start from the maximally mixed
one-excitation input and apply the retained heat contraction $A_J^{1/2}$.
The success probability is $Z_J/D$, bounded below at fixed accuracy
because both $Z$ and $D$ scale as $h^{-1}$. An optical unitary dilation
realizes the contraction; universal interferometer synthesis
\cite{Clements2016} costs $O(D^2)$ elements and depth $O(D)$.
These are costs of a programmed state source with known matrix data.
Input heralding, rejected trials, mode-dependent loss, and reset remain
part of its resource budget. Balanced direct eigenvector preparation
offers a different tradeoff, described in
Appendix~\ref{sgt:prep:appendix}.

For $\gamma=1/2$, subordination gives an exact alternative using a
no-loss evolution generated by $L$:
\begin{equation}
 \ee^{-\beta\sqrt L}=\int_0^\infty
 \frac{\beta\ee^{-\beta^2/(4u)}}{2\sqrt\pi u^{3/2}}
 \ee^{-uL}\,\dd u .
 \label{sgt:main:subordination}
\end{equation}
The successful, unnormalized outputs must be pooled before
normalization. The sampling law has infinite mean time, but a controlled
cutoff gives fixed accuracy; alternatively, aborting at the first
detected loss gives mean attempted evolution time $O(h)$ in the
dimensionless generator units. Bounding the physical decay rate rescales
this time to $O(1)$ at fixed accuracy. The construction thus pays for
the growing generator bandwidth instead of interpreting a shrinking
mathematical time as unlimited laboratory speed. A local bulk stiffness
matrix reduces to a generally dense boundary operator, so its static
Schur complement alone does not implement the required dynamics.

\section{A full Bose ensemble is a different resource}
\label{sgt:main:bose}

Now specify a common number-conserving Fock model on
$\mathcal F_s(\mathcal H)$:
\begin{equation}
\begin{gathered}
 \mathbb H_h=\dd\Gamma(H_h),\\
 \varrho_h=\det(I-A_h)\,\Gamma(A_h),\\
 N_h=A_h(I-A_h)^{-1}.
 \label{sgt:main:fock}
\end{gathered}
\end{equation}
The chemical potential is zero and the reference has its thermal
state on $P$ and exact vacuum on $Q$. This ensemble has variable
particle number; it is not the canonical ensemble at a fixed number
of excitations. Energies are excitation energies relative to the
declared vacuum.

For an orthogonal detected subspace $E$, efficiency $0\le\eta\le1$,
and $0\le z\le1$, the occupation generating function is
\begin{equation}
 \mathbb E z^K=\det\bigl(I+\eta(1-z)EN_hE\bigr)^{-1}.
 \label{sgt:main:PGF}
\end{equation}
This familiar bosonic determinant \cite{Klich2002,Quesada2018} does not require
$E$ to commute with $A_h$. In particular, an ideal $Q$ receiver has
zero false-positive probability and
\begin{equation}
 p_{\rm off}\le
 \exp\{-\log(1+\eta)[C_Q(\tau)+o(1)]/h\}.
 \label{sgt:main:off}
\end{equation}
This is distinguishability of an entire many-excitation ensemble.
The same determinant mechanism underlies established thermal
source-discrimination receivers \cite{Lu2018}.

The total resource follows from the continuum law. Define
\begin{align}
 \mathcal L(\tau)&=\sum_{k\ge1}C(k\tau)/k,&
 \mathcal E(\tau)&=\sum_{k\ge1}C_1(k\tau)/k,\nonumber\\
 \mathcal V(\tau)&=\sum_{k\ge1}C_2(k\tau)/k,&
 \mathcal B(\tau)&=\sum_{k\ge1}C(k\tau).
 \label{sgt:main:bosecoeff}
\end{align}
For every $\gamma>0$, $h\log\mathcal Z_B\to\mathcal L$,
$h\beta\langle\mathbb H_h\rangle\to\mathcal E$, and
$h\beta^2\operatorname{Var}(\mathbb H_h)\to\mathcal V$.
The Fock entropy satisfies $hS(\varrho_h)\to\mathcal L+\mathcal E$.
For $0<\gamma<1$, the number sum also converges and
$h\Tr N_h\to\mathcal B$. At $\gamma=1/2$, mean particle number is
$\Theta(h^{-1})$ and mean excitation energy is $\Theta(h^{-3/2})$.
Agreement in a subsequent high-temperature limit concerns ratios of
these extensive coefficients, not vanishing additive Fock-entropy
differences. Indeed the strict finite-$\tau$ partition deficit implies
an exponentially small ratio of the split and collapsed Bose partition
functions.

The probability that a full Bose sample contains exactly one excitation
is $Z_h/\mathcal Z_B\le Z_h\ee^{-Z_h}$, exponentially small in $h^{-1}$.
Controlled preparation avoids that particular postselection cost;
unconditioned Bose readout instead uses the many available excitations.
These are two consistent experiments, not interchangeable copy counts.

\section{Robust counting and the role of infrared rejection}
\label{sgt:main:robust}

A dark-port rule is sensitive to a false-click floor. With a sorter $V$
and reference covariance supported on $P$, leakage amplitude
$\theta=\norm{QVP}$ gives reference population at most
$\theta^2\Tr N_{0,h}$. Thus fixed crosstalk can make an any-click rule
unreliable as the number of occupied modes grows. This is the familiar
reason to replace a dark-port decision by a calibrated count threshold
\cite{Schlichtholz2024}.

For the present geometry, let $0<\gamma<1$ and let a common detection
effect differ from $Q$ by a sufficiently small fixed operator norm.
The two mean counts still differ by a positive multiple of $h^{-1}$.
The uniform estimate~\eqref{sgt:main:weakS1} bounds the largest Bose
occupation by $O(h^{-\gamma})$. A geometric-count moment-generating
bound then gives a midpoint-threshold error
\begin{equation}
 P_{\rm err}\le \exp[-c h^{\gamma-1}].
 \label{sgt:main:threshold}
\end{equation}
Independent calibrated Poisson background with mean $O(h^{-1})$
is allowed. The detector resolves total occupation, with a dynamic
range of order $h^{-1}$; a single bucket-detector bit is a different
measurement.

The loss in the exponent comes from very highly occupied infrared
modes. A common filter can remove them under both hypotheses. Choose
a small fixed $b>0$ and exclude a $P$-invariant span containing both
low-energy spaces below $b/h$. Its rank is $O(b/h)$ and its lost mean
population is at most $O(b^{1-\gamma}/h)$. For small $b$, a positive
transverse mean gap remains. A subsequent common finite Fourier
restriction produces a projection $\Xi_h$ of rank $O(h^{-1})$,
commuting with $P$, on which both covariances are uniformly bounded.
A fixed sufficiently small error within this retained space therefore
permits
\begin{equation}
 P_{\rm err}\le\exp(-c/h).
 \label{sgt:main:filtered}
\end{equation}
The filter is constructed from the two known candidate operators and
is used unchanged under both hypotheses.

This conclusion has a quantitative stopband condition. For a physical
passive contraction $T_h$ into $O(h^{-1})$ output modes, set
$\ell_h=\norm{T_h(I-\Xi_h)}$ and assume
$\norm{\Xi_hT_h^*T_h\Xi_h-\Xi_hQ}\le e$.
For fixed positive efficiency, if $e+2\ell_h+\ell_h^2$ stays below
the fixed tolerance preserving the positive mean gap, then
\begin{equation}
 P_{\rm err}\le
 \exp\left[-\frac{c}{h(1+\ell_h^2h^{-\gamma})}\right].
 \label{sgt:main:leakage}
\end{equation}
In particular, amplitude leakage $O(h^{\gamma/2})$ retains the full
exponent; at $\gamma=1/2$ this is $O(h^{1/4})$. The distinction between
fixed errors inside the passband and leakage from a highly populated
stopband is essential. Appendix~\ref{sgt:bose:appendix} proves the
finite-space construction, count concentration, and leakage estimate.

\section{An exact finite circular instance}
\label{sgt:main:finite}

A rotationally invariant example provides finite numbers without a
Weyl approximation. In the unit Dirichlet disk take circles of radii
$99/200$ and $101/200$, each carrying mass $1/2$ uniformly in angle.
The reference is the mass-one circle of radius $1/2$, aligned with
the symmetric layer mode. For angular order $n\ge1$ the exact inverse
block is
\begin{equation}
 (R_n)_{ij}=\frac1{8\pi n}
 \left[\left(\frac{\min(r_i,r_j)}{\max(r_i,r_j)}\right)^n
                         -(r_ir_j)^n\right].
 \label{sgt:main:circle}
\end{equation}
Both angular signs are included. Use $\gamma=1/2$, $\beta=1/20$,
and measure the antisymmetric-like output in $50\le|n|\le500$.
Allow a separate fixed real sorting-angle error in
$[-1/10,1/10]$ radians in each sector, collection efficiency $7/10$,
and independent Poisson dark counts of total mean at most $1/10$.
Declare the split geometry when the total recorded occupation is at
least five. Validated arithmetic gives the bounds in Table~\ref{sgt:main:table}.
\begin{table}[b]
 \caption{Rigorous bounds for the finite circular count experiment.
 They bound the specified receiver error, not the optimal error.}
 \label{sgt:main:table}
 \begin{ruledtabular}
 \begin{tabular}{lc}
 Quantity & Upper bound\\
 \hline
 False positive & $0.0057$\\
 False negative & $0.00072$\\
 Equal-prior error & $0.0032$
 \end{tabular}
 \end{ruledtabular}
\end{table}
There are $902$ measured output modes and $1804$ prepared-band modes.
Mean prepared occupations are approximately $61.7931$ and $81.3157$
for the split and reference states. The corresponding excitation
energies are approximately $3550.20$ and $3590.81$ in the stated units.
The bound is uniform over the specified sectorwise angles and dark-count
means. Its derivation uses exact $2\times2$ functional calculus,
geometric occupation generating functions, and interval-certified
Chernoff bounds; the complete calculation is in
Appendix~\ref{sgt:finite:certificate}.

The entire two-layer angular band is a reducing subspace. Its Gaussian
marginal may be prepared alone while reproducing these measured
statistics; omitted modes of a full ensemble are traced out. The finite
instance assumes ideal rejection outside the selected angular band.
It establishes neither matched brightness nor finite-parameter thermal
scalar agreement. The asymptotic leakage law in
Sec.~\ref{sgt:main:robust} is a separate result, not an uncomputed
finite tolerance for this example.

\section{Physical realization and interpretation}
\label{sgt:main:physics}

Two implementation facts determine how to interpret these results.
First, a finite gauge-invariant Gaussian state with covariance $N$ has
a positive coherent-state representation: sample a complex Gaussian
amplitude with covariance $N$ and prepare that coherent amplitude.
This established representation \cite{Lu2018} gives a programmed
optical implementation of any declared finite marginal. For a
general detector subspace the correct input is its compressed
covariance, not the Bose function of a compressed Hamiltonian. Dense
covariance factorization and sampling cost $O(D^3)$ once and $O(D^2)$
per draw, respectively, with $D$ mode controls. This simulator knows
the covariance; an unknown-specimen sensor additionally needs a
physical coupling that supplies it.

Second, a common one-particle alignment does not automatically align
mechanical vacua \cite{BruneauDerezinski2007}. In fixed annihilation
coordinates a real symmetric
positive matrix $H$
gives the number-conserving Hamiltonian
\begin{equation}
 \dd\Gamma(H)=\tfrac12(x^{\mathsf T}Hx+p^{\mathsf T}Hp-\Tr H).
 \label{sgt:main:fixedvac}
\end{equation}
An ordinary oscillator instead has
$(p^2+\omega^2x^2)/2$ and frequency-dependent annihilation operators.
Relative to a reference frequency, its vacuum is squeezed, with
$r=\tfrac12\log(\omega/\omega_{\rm ref})$ and mean reference occupation
$\sinh^2r$. Raising a mechanical frequency is therefore not the same
operation as making a fixed optical port empty. A native realization
must pay for canonical conversion, added noise, bandwidth, and reset,
as in the broader optomechanical setting \cite{Aspelmeyer2014}.
For example, restoring a fixed frequency unit $\Omega_*$ at
$\gamma=1/2$ gives $\beta=\hbar\Omega_*/(k_BT)=\tau\sqrt h$.
The joint regime then has $T\propto h^{-1/2}$ and a retained frequency
band of order $\Omega_*h^{-1/2}$ at fixed truncation accuracy.
An optical simulator at a fixed carrier reproduces the programmed
occupation statistics, not these mechanical temperature and energy units.

The mathematical infinite-energy reference has controlled finite
replacements. In a one-excitation model, retain longitudinal dispersion
and add a gap $\Delta$ to each of the $M-1$ transverse channels. Their
total probability is
\begin{equation}
 \frac{(M-1)\ee^{-\beta\Delta}}
             {1+(M-1)\ee^{-\beta\Delta}}.
 \label{sgt:main:gapone}
\end{equation}
For a full Bose reference with $D_Q$ detected transverse modes,
their expected detected occupation is bounded by
$\eta D_Q/(\ee^{\beta\Delta}-1)$. The condition
$\beta\Delta-\log D_Q\to+\infty$ suffices for an any-click vacuum
approximation, whereas a calibrated
count threshold only requires an adequate mean margin. A scalar gap
alone on infinitely many modes has infinite thermal trace; longitudinal
confinement must remain present.

A collapsing geometry can
retain transverse information after a broad set of normalized thermal
summaries has converged. Accessing it can use a controlled excitation
or a many-excitation ensemble, with different costs. In the latter,
the largest modal occupations, rather than the total signal alone,
govern robustness. Common infrared rejection separates the signal
budget from the fluctuation budget and yields an explicit interpolation
between the two error exponents. The finite circular certificate gives
a directly reproducible target for a common-mode emulator. A native
geometric device would further require a uniform thin-layer, carrier,
and mode-conversion approximation over its occupied band.

\begin{acknowledgments}
The author declares the substantial use of consumer-grade Gemini-3 and GPT-6 models for literature search, exploring and developing ideas, codes and proofs, drafting and LaTeX typesetting; the author takes full and sole responsibility for all contents of this work.
\end{acknowledgments}

\section*{Data availability}
Numerical results reported here are reproducible from the verification
scripts and certificate data in the repository - Repository - 
\href{https://github.com/arun-chandru/colliding_layers}{github.com/arun-chandru/colliding\_layers}.


\appendix
% Standalone mathematical content, included after the manuscript's \appendix.
% Uses only standard AMS commands and the shared macros Tr, dd, ee, id, norm.
% All hypotheses and all auxiliary arguments used below are stated here.

\section{Physical mass, alignment, and the complete joint trace law}
\label{sgt:geo:setup}
\label{sgt:geo:appendix}

Let $\Omega\subset\mathbb R^2$ be a fixed bounded open set, and let
$\Gamma$ be a smooth embedded closed curve of length $\ell$ (the
quantity denoted by $L$ in the main text).  Write
$\mathbf r(s)$ for its arclength parametrization, with
$s\in\mathbb R/\ell\mathbb Z$, and $n(s)$ for a smooth unit normal.
A closed tubular neighborhood of $\Gamma$ is assumed to lie compactly
inside $\Omega$.  This allows any fixed absorbing masks a positive
distance from the tube, without regularity assumptions on their
distant boundaries.  Fix a smooth positive function $w$ with
$\int_\Gamma w\,\dd s=1$, distinct fixed heights
$a_1<\cdots<a_M$, and positive fixed weights $\alpha_i$ with
$\sum_i\alpha_i=1$.  For sufficiently small $h>0$ put
\begin{equation}
 T_{i,h}(s)=\mathbf r(s)+ha_i n(s),\qquad
 \nu_h=\sum_{i=1}^M\alpha_i(T_{i,h})_*(w\,\dd s).
 \label{sgt:geo:mass}
\end{equation}
The sheets are then disjoint and embedded.  The measure in
\eqref{sgt:geo:mass} is the physical mass: its expression in the
coordinate $s$ is $\alpha_iw(s)\dd s$.  In particular there is no
arclength Jacobian to insert when a sheet is moved.

For clarity, the Green inverse used here has a direct variational
definition.  Give $H^1_0(\Omega)$ the energy inner product
$\int_\Omega\nabla u\cdot\overline{\nabla v}\,\dd x$, with no factor
$1/2$, and let $J_hu$ be its traces on the sheets, considered as an
element of $L^2(\nu_h)$.  The usual local trace estimate in a smooth
tube, followed by Poincare's inequality on the bounded carrier, makes
$J_h$ bounded.  The trace lies locally in $H^{1/2}$; compactness of
$H^{1/2}(\Gamma)\hookrightarrow L^2(\Gamma)$ makes $J_h$ compact.
Its range is dense: smooth data on finitely many disjoint sheets can
be extended separately inside disjoint smaller tubes.  The physical
closed form is
\[
\begin{gathered}
 \mathcal E_h(f)=\inf\left\{
   \int_\Omega|\nabla u|^2\,\dd x:J_hu=f\right\},
 \\
 f\in J_hH^1_0(\Omega).
\end{gathered}
\]
The infimum is attained on the energy-orthogonal complement of
$\ker J_h$.  That complement is complete in energy norm; the bounded
trace map identifies it with this form domain, proving closedness.
The inverse of its positive self-adjoint generator is $J_hJ_h^*$.
Indeed $u=J_h^*f$ is the unique weak solution with sheet source $f\nu_h$,
and its trace is $J_hJ_h^*f$.  Density of the range proves injectivity of
this positive compact inverse.  Denote the generator by $L_h$ and its
inverse by $R_h=L_h^{-1}$.

Let $G_\Omega$ be the Dirichlet Green kernel for this energy.  The weak
solution representation gives the integral kernel $G_\Omega(x,y)$ on
$L^2(\nu_h)$.  Conjugating by the exact unitary map
\begin{equation}
\begin{gathered}
 (U_hf)_i(s)=\sqrt{\alpha_iw(s)}\,f(T_{i,h}(s))
 \\
\hbox{onto}\\
 \mathcal H=L^2(\Gamma,\dd s)\otimes\mathbb C^M
 \label{sgt:geo:unitary}
\end{gathered}
\end{equation}
gives the matrix kernel
$\sqrt{w(s)w(t)\alpha_i\alpha_j}\,
G_\Omega(T_{i,h}(s),T_{j,h}(t))$, again denoted by $R_h$.
All logarithmic diagonal singularities are square integrable, so
$R_h\in\mathfrak S_2$.  Thus all positive-time heat traces and their
fixed polynomial energy moments below are finite: the decreasing
inverse eigenvalues satisfy $r_j\leq\norm{R_h}_2/\sqrt j$, which already
implies summability of $\exp[-t r_j^{-\gamma}]$ for every $t,\gamma>0$.
The stronger estimate used for Bose singularities elsewhere is not an
input to the argument in this appendix.

Set
\[
\begin{gathered}
 v=(\sqrt{\alpha_1},\ldots,\sqrt{\alpha_M})^t,\\
 P=vv^*,\\
 Q=I_M-P,\\
 z(s)=\sqrt{w(s)}.
\end{gathered}
\]
The corresponding projections on $\mathcal H$ have the same names.
The collapsed physical space is embedded as $L^2(\dd s)\otimes v$.
If $R$ is the base-curve Green inverse conjugated from $L^2(w\dd s)$
to $L^2(\dd s)$, its aligned extension is $R_0=R\otimes P$.
The orthogonal $Q$ space has inverse zero and heat zero.  It represents
an alignment complement, not physical modes of zero energy.

For $q>0$ define
\begin{equation}
\begin{gathered}
 T(q)_{ij}=\sqrt{\alpha_i\alpha_j}\,\ee^{-q|a_i-a_j|},
 \\
 B(s,\xi)=\frac{w(s)}{2|\xi|}T(|\xi|),\\
 B_0(s,\xi)=\frac{w(s)}{2|\xi|}P.
 \label{sgt:geo:symbol}
\end{gathered}
\end{equation}
Fix $\gamma>0$ and $\tau>0$.  With the value at zero defined to be zero,
write
\[
\begin{gathered}
 f_{\tau,\gamma}(b)=\ee^{-\tau b^{-\gamma}},\\
 X_h=R_h/h,\\
 X_{0,h}=R_0/h,\\
 A_h(\tau)=f_{\tau,\gamma}(X_h),\\
 A_{0,h}(\tau)=f_{\tau,\gamma}(X_{0,h}).
\end{gathered}
\]
These are the heat operators at inverse temperature
$\beta=\tau h^\gamma$ for $L_h^\gamma$.  In particular $\gamma=1/2$
weights single-excitation frequencies if $L_h$ gives squared
frequencies.  Thus $H_h=L_h^\gamma$ is the one-particle Hamiltonian,
and $A_h(\tau)=\exp(-\beta H_h)$ is its heat operator.  Let
$Z_h=\Tr A_h$, $Z_{0,h}=\Tr A_{0,h}$ and
$\rho_h=A_h/Z_h$, $\rho_{0,h}=A_{0,h}/Z_{0,h}$.  Put
\begin{gather}
 F(s,\xi)=f_{\tau,\gamma}(B(s,\xi)),\\
 F_0(s,\xi)=\ee^{-\tau(2|\xi|/w(s))^\gamma}P,\notag\\
 C_\gamma(\tau)=\frac1{2\pi}\int_\Gamma\int_{\mathbb R}
                         \operatorname{tr}F(s,\xi)\,\dd\xi\dd s,\\
 G_\gamma=\Gamma(1+1/\gamma),\notag\\
 C_{\mathrm{col},\gamma}(\tau)=
             \frac{G_\gamma}{2\pi\tau^{1/\gamma}}.
 \label{sgt:geo:coefficients}
\end{gather}
Lowercase $\operatorname{tr}$ denotes a finite matrix trace.
Values at $\xi=0$ have no effect on these integrals.

The complete actual-operator conclusions are
\begin{gather}
 hZ_h\longrightarrow C_\gamma(\tau),\\
 hZ_{0,h}\longrightarrow C_{\mathrm{col},\gamma}(\tau),
 \label{sgt:geo:heat-law}\\
 \frac12\norm{\rho_h-\rho_{0,h}}_1
 \longrightarrow
\\
 D_\gamma(\tau):=\frac1{4\pi}\int_\Gamma\int_{\mathbb R}
\\
 \operatorname{tr}\left|
 \frac{F}{C_\gamma(\tau)}-
 \frac{F_0}{C_{\mathrm{col},\gamma}(\tau)}
 \right|\,\dd\xi\dd s,
 \label{sgt:geo:distance}\\
 \Tr(Q\rho_h)\longrightarrow
 p_{Q,\gamma}(\tau):=\frac{C_{Q,\gamma}(\tau)}{C_\gamma(\tau)},\\
 C_{Q,\gamma}(\tau)=\frac1{2\pi}\int_\Gamma\int_{\mathbb R}
                    \operatorname{tr}(QF)\,\dd\xi\dd s.
 \label{sgt:geo:q-law}
\end{gather}
For $M\geq2$ these coefficients obey
$0<p_{Q,\gamma}(\tau)\leq D_\gamma(\tau)<1$.
All data $M,a_i,\alpha_i,\Gamma,w,\Omega,\gamma,\tau$ are fixed before
$h\downarrow0$.  Constants may depend on all of them.  Endpoint limits
in $\tau$ below are subsequent limits of these functions.  No uniform
claim is made when $M$ grows, a weight or height separation vanishes,
a mask approaches the tube, or $\tau$ is chosen as an arbitrary function
of $h$.

\subsection{A positive-resolvent comparison that retains every mode}
\label{sgt:geo:resolvent}

We prove a slightly larger functional-calculus statement.  Let $g$ be
real and continuous on $[0,\infty]$, with $g(0)=0$ and a finite limit
of $g(x)/x^2$ at zero.  For $\mathcal T(A)=A(I+A)^{-1}$ and
$\Phi(y)=g(y/(1-y))$, the function $\Phi(y)/y^2$ is continuous on
$[0,1]$.  For every $\eta>0$ a polynomial $p$ therefore satisfies
\begin{equation}
 |\Phi(y)-y^2p(y)|\leq\eta y^2\quad(0\leq y\leq1).
 \label{sgt:geo:polynomial}
\end{equation}
Suppose $U=\mathcal T(A)$, $V=\mathcal T(D)$ for positive compact operators, with
$\Tr U^2+\Tr V^2\leq C/h$, and
$U-V=J+E$, where $\operatorname{rank}J\leq r$,
$\norm{J}\leq C_J$ and $\norm{E}_2\leq e$.
For a monomial of degree $k\geq2$ use
\[
 U^k-V^k=\sum_{j=0}^{k-1}U^j(J+E)V^{k-1-j}.
\]
Each $J$ term has trace norm at most $C_Jr$.  In each $E$ term at
least one outside exponent is positive, and the corresponding factor
has Hilbert--Schmidt norm at most
$\max(\norm{U}_2,\norm{V}_2)$.  Schatten Holder's inequality then gives
\begin{equation}
 h\norm{g(A)-g(D)}_1
 \leq C\eta+C_p\bigl(hC_Jr+\sqrt h\,e\bigr).
 \label{sgt:geo:comparison}
\end{equation}
Here the two approximation errors in \eqref{sgt:geo:polynomial} cost
at most $\eta\Tr(U^2+V^2)$, by spectral calculus.

For a decomposition $A-D=J_{\rm raw}+E_{\rm raw}$, the exact identity
\begin{equation}
 \mathcal T(A)-\mathcal T(D)=(I+A)^{-1}(A-D)(I+D)^{-1}
 \label{sgt:geo:resolvent-identity}
\end{equation}
preserves the rank bound and does not increase the residual
Hilbert--Schmidt norm.  If $\norm{E_{\rm raw}}\leq C$ uniformly, the
transformed finite-rank term has norm at most $1+C$, since the
difference of two positive contractions has norm at most one.
Even without that operator-norm bound, a uniformly bounded rank-$r$
decomposition is available.  Set $K=U-V$ and keep its $r$ largest
singular values, with their signs, to obtain $J$.  Then $\norm{J}\leq1$.
The squared Hilbert--Schmidt residual is the sum of the remaining
squared singular values, and is no larger than for any other rank-$r$
approximation; applying this to the rank-$r$ term already supplied by
the resolvent identity gives $\norm{K-J}_2\leq\norm{E_{\rm raw}}_2$.
For completeness, that minimizing property follows by projecting onto
the range of any proposed rank-$r$ approximation: its discarded
Hilbert--Schmidt norm is at least the sum after the largest $r$
eigenvalues of $KK^*$, by the variational principle for their partial
sum.  Spectral truncation attains the bound.  The difference $K$ is
Hilbert--Schmidt already, because the resolvent identity writes it as
a finite-rank term plus such a residual.  Thus one may always take
$C_J=1$ after this additional truncation; no bound on the size of
$J_{\rm raw}$ is needed.  If only the budget for $V$ is initially known,
\[
 \norm{U}_2\leq\norm{V}_2+C_J\sqrt r+e
\]
supplies it whenever the displayed right side is $O(h^{-1/2})$.
In particular $hr=o(1)$ and $\sqrt h\,e=o(1)$ imply a whole-operator
$o(h^{-1})$ trace-norm comparison.  This proof uses no scalar
Lipschitz assertion in trace norm and no short-time extrapolation of
a fixed-temperature bound.

\subsection{The curved kernel and a manifestly positive model}
\label{sgt:geo:kernel}

Write $k_n=2\pi n/\ell=\varrho n$.  Define the positive matrix
convolution operator $\mathcal C_h$ by its Fourier multipliers
\begin{equation}
\begin{gathered}
 \mathcal C_h(0)=I_M,\\
 \mathcal C_h(n)=\frac{T(h|k_n|)}{2|k_n|}\\
(n\ne0),
 \\
 \mathcal A_h=h^{-1}M_z\mathcal C_hM_z.
 \label{sgt:geo:model}
\end{gathered}
\end{equation}
For the collapsed model use the scalar convolution $\mathcal C_0$
with multipliers $\mathcal C_0(0)=1$,
$\mathcal C_0(n)=(2|k_n|)^{-1}$ for $n\ne0$, and
$\mathcal A_{0,h}=h^{-1}(M_z\mathcal C_0M_z)\otimes P$.
The finitely many zero-frequency choices affect no leading trace.
Positivity follows from
\[
 \ee^{-q|a-b|}=\frac1\pi\int_{\mathbb R}
       \frac{q\ee^{it(a-b)}}{q^2+t^2}\,\dd t.
\]
A nonzero finite exponential sum at distinct heights is not identically
zero, so the associated quadratic form is strictly positive for $q>0$.
Consequently $T(q)>0$, $\operatorname{tr}T(q)=1$, and $\norm{T(q)}\leq1$.
At most $M$ eigenvalues occur at each Fourier index and each is
$O((1+|n|)^{-1})$.  Min--max after multiplication by bounded $z$ yields
\begin{equation}
\begin{gathered}
 \Tr\mathcal T(\mathcal A_h)^2+\Tr\mathcal T(\mathcal A_{0,h})^2
\\
 \leq C\sum_{j\geq1}\min\{1,(hj)^{-2}\}\leq C'/h.
 \label{sgt:geo:model-budget}
\end{gathered}
\end{equation}

The regular part
$G_\Omega(x,y)+(2\pi)^{-1}\log|x-y|$ is smooth on the product of
the closed tube with itself.  Both partial Laplacians vanish
distributionally there; their sum is elliptic in four variables, so
interior regularity includes the diagonal.  This argument is local
and does not require smooth distant boundary components.  Define
\[
\begin{gathered}
 \ell_*(u)=\frac\ell\pi\left|\sin\frac{\pi u}{\ell}\right|,
 \\
 \Lambda_\delta(u)=-\frac1{4\pi}
              \log\left(1+\frac{\delta^2}{\ell_*(u)^2}\right),
 \\
 D(s,t)=n(s)\cdot\nabla_xG_\Omega(\mathbf r(s),\mathbf r(t)).
\end{gathered}
\]
The latter off-diagonal expression extends smoothly: the normal
numerator in the logarithmic part vanishes to order $(s-t)^2$, as
does the chord-length squared denominator.  The ratio of the geometric
chord to $\ell_*(s-t)$ is smooth and positive across the diagonal.
It follows that $R=M_z\mathcal C_0M_z+\mathcal S_0$, with
$\mathcal S_0$ a smooth-kernel operator.

For $\delta=h|a_i-a_j|$, the required uniform kernel expansion is
\begin{gather}
 G_\Omega(T_{i,h}(s),T_{j,h}(t))
\\
 =G_\Omega(\mathbf r(s),\mathbf r(t))+\Lambda_\delta(s-t)
\\
     +h\{a_iD(s,t)+a_jD(t,s)\}+\mathcal E_{h,ij}(s,t),\notag\\
 |\mathcal E_{h,ij}(s,t)|
 \leq Ch^2+Ch\frac{\delta^2}{u^2+\delta^2},
 \label{sgt:geo:strip}
\end{gather}
where $u$ is the short periodic coordinate, and the fraction is zero
when $\delta=0$.  Here and below kernels may be specified off their
null diagonal.  To verify the singular part, put
$p=\mathbf r(s)-\mathbf r(t)$,
$b=n(s)-n(t)$, $c=n(s)+n(t)$,
$m=(a_i+a_j)/2$, and $d=(a_i-a_j)/2$.
The changed chord is $p+h(mb+dc)$ and $b\cdot c=0$.
The smooth quotients satisfy
\[
\begin{gathered}
 \frac{p\cdot b}{|p|^2}=O(1),\\
 \frac{p\cdot c}{|p|^2}=O(|u|),\\
 \frac{|b|^2}{|p|^2}=O(1),\\
 \frac{|c|^2}{|p|^2}=\frac4{\ell_*(u)^2}+O(1).
\end{gathered}
\]
These follow by Taylor expansion of an arclength parametrization and
$n\perp\mathbf r'$.  Expand the logarithm of
$|p+h(mb+dc)|^2/|p|^2$ about $1+\delta^2/\ell_*(u)^2$.
The relative perturbation is $O(h)$, its quadratic logarithm error is
$O(h^2)$, and replacing the denominator of the linear term by one
costs at most the second term in \eqref{sgt:geo:strip}.
The smooth regular part and points away from the diagonal have the
ordinary $O(h^2)$ Taylor remainder.  This proves the estimate, including
unequal, uncentered profiles.  A uniformly $O(h^2)$ pointwise remainder
would in general be false.  Squaring and integrating its strip term
gives $O(h^2\delta)=O(h^3)$; hence the weighted matrix remainder has
Hilbert--Schmidt norm $O(h^{3/2})$.

The exact Fourier multipliers of $\Lambda_\delta$, for convolution
against $\dd s$, are
\begin{equation}
\begin{gathered}
 b_\delta=2\operatorname{arsinh}(\varrho\delta/2),\\
 \Lambda_\delta(0)=-\frac{b_\delta}{2\varrho},\\
 \Lambda_\delta(n)=-\frac{1-\ee^{-b_\delta|n|}}
                            {2\varrho|n|}\\
(n\ne0).
 \label{sgt:geo:log-fourier}
\end{gathered}
\end{equation}
For example factor
$\sin^2(\theta/2)+\sinh^2(b_\delta/2)
=(\cosh b_\delta-\cos\theta)/2$ and expand the logarithm in its
integrable Fourier series.  Its constant term is $b_\delta$, and its
cosine coefficient is $2(1-\ee^{-b_\delta n})/n$, giving precisely
\eqref{sgt:geo:log-fourier}.  For each fixed nonzero height difference,
\[
\begin{gathered}
 |b_{h|a_i-a_j|}-\varrho h|a_i-a_j||\leq Ch^3,\\
 \min\{b_{h|a_i-a_j|},\varrho h|a_i-a_j|\}\geq ch.
\end{gathered}
\]
The mean value theorem therefore bounds the discrepancy of nonzero
Fourier multipliers by $Ch^3\ee^{-ch|n|}$.  Its squared sum is
$O(h^5)$; diagonal matrix entries have zero discrepancy.
One must use the positive model \eqref{sgt:geo:model}, not assume
positivity of the matrix obtained from the exact periodic replacement.
Collecting terms gives
\begin{equation}
\begin{gathered}
 R_h=M_z\mathcal C_hM_z+\mathcal S_0\otimes P+h\mathcal S_1
       +J_{0,h}+\mathcal R_h+\mathcal E_h,
 \\
 \operatorname{rank}J_{0,h}\leq M,\\
 \norm{\mathcal R_h}_2=O(h^{5/2}),\\
 \norm{\mathcal E_h}_2=O(h^{3/2}).
 \label{sgt:geo:normal-form}
\end{gathered}
\end{equation}
Here $\mathcal S_1$ is the fixed smooth matrix kernel formed from the
normal derivatives.  The zero-frequency discrepancy, sandwiched by
$z$, has range in $\operatorname{span}\{z\}\otimes\mathbb C^M$, giving
the stated rank.  Its norm will not be needed.

If $\Pi_N$ projects onto $|n|\leq N$ in every component, any fixed
smooth matrix kernel $\mathcal S$ satisfies
$\norm{\mathcal S-\Pi_N\mathcal S\Pi_N}_2\leq C_kN^{-k}$ for every
integer $k$, by repeated integration by parts in both Fourier indices.
Take $N=\lceil h^{-1/2}\rceil$, $k\geq4$, and divide
\eqref{sgt:geo:normal-form} by $h$.  The result is a finite-rank term
of rank $O(h^{-1/2})$ and a Hilbert--Schmidt residual $O(h^{1/2})$.
Equations \eqref{sgt:geo:resolvent-identity} and
\eqref{sgt:geo:model-budget} first give the actual resolvent-square
budget and then \eqref{sgt:geo:comparison} gives
\begin{equation}
\begin{gathered}
 h\norm{g(X_h)-g(\mathcal A_h)}_1\longrightarrow0,
 \\
 h\norm{g(X_{0,h})-g(\mathcal A_{0,h})}_1\longrightarrow0.
 \label{sgt:geo:actual-model}
\end{gathered}
\end{equation}
Both operators compared are positive.  No positivity is attributed to
an intermediate truncated geometric perturbation.

\subsection{Matrix Fourier words, low and high frequencies, and sharp bands}
\label{sgt:geo:fourier}

Choose $0<2d<R$ and a smooth $0\leq\chi\leq1$ supported where
$d\leq|\xi|\leq2R$, equal to one where $2d\leq|\xi|\leq R$.
Replace the nonzero multipliers of $\mathcal C_h/h$ by
$\chi(hk_n)T(|hk_n|)/(2|hk_n|)$, and set the zero multiplier to zero.
Keep the two multiplications by $z$; call the result
$\mathcal A_h^{d,R}$.  Use the same construction for the reference.
These positive operators have rank $O(R/h)$ and uniformly bounded
norm at fixed $d,R$.  They are below their full models in operator
order, so eigenvalue comparison preserves the resolvent-square trace
budget.  The discarded low-frequency part has rank $O(d/h)+O(1)$.
The high-frequency part obeys
\begin{equation}
 \norm{E_{\rm hi}}\leq C/R,\qquad
 \norm{E_{\rm hi}}_2^2\leq C/(hR).
 \label{sgt:geo:frequency-tail}
\end{equation}
Indeed these follow by summing $(h|n|)^{-2}$ over $|hk_n|>R$ and
using $\norm{T(q)}\leq1$; bounded multiplication by $z$ changes only
constants.  After the resolvent identity the finite-rank part has
norm at most $1+C/R$, independently of $d$.  Consequently
\begin{equation}
 \limsup_{h\downarrow0}h\norm{g(\mathcal A_h)
                              -g(\mathcal A_h^{d,R})}_1
 \leq C\eta+C_p(d+R^{-1/2}),
 \label{sgt:geo:cutoff-removal}
\end{equation}
and identically for the reference.  For each fixed approximation
polynomial first let $d\downarrow0$, $R\uparrow\infty$, and then
$\eta\downarrow0$.  This controls all original Fourier modes outside
the annulus; no arbitrary eigenvalue reindexing is involved.

Here is the trace calculation on the fixed annulus.  Expand a word
in operators $M_{z_j}m_j(hD_s)M_{z_j}$, with smooth scalar $z_j$ and
smooth compactly supported matrix multipliers $m_j$, in the basis
$\ell^{-1/2}\ee^{ik_ns}$.  Each multiplication shifts the index by an
integer and has rapidly decreasing, absolutely summable Fourier
coefficients.  For any fixed list of shifts $r_j$,
$m_j(hk_{n+r_j})-m_j(hk_n)=o(1)$ uniformly in $n$.
The number of contributing $n$ is $O(h^{-1})$ uniformly in the shifts,
by anchoring to any one compactly supported multiplier factor; the
remaining factors are bounded.  The absolute sum of the products of
Fourier coefficients is finite.  Dominated convergence therefore
permits summing over the shift lists.  The zero-total-shift condition
reconstructs the spatial average of the position factors.  The
remaining frequency sum is a Riemann sum of spacing $2\pi h/\ell$.
Thus, retaining matrix order,
\begin{equation}
\begin{gathered}
 h\Tr(\mathcal B_{1,h}\cdots\mathcal B_{r,h})
\\
 \longrightarrow\frac1{2\pi}\int_\Gamma\int_{\mathbb R}
        \operatorname{tr}(b_1(s,\xi)\cdots b_r(s,\xi))
                         \,\dd\xi\dd s.
 \label{sgt:geo:word-law}
\end{gathered}
\end{equation}
The same calculation permits smooth scalar position insertions,
fixed matrix insertions, and a smooth bounded Fourier multiplier
$\psi(hD_s)$: its shifted values also converge uniformly for any
fixed shift on the relevant compact frequency set.  No factors have
been commuted in this calculation.

Uniform polynomial approximation on the bounded spectra of the
cutoff models, with approximating polynomials zero at zero, extends
\eqref{sgt:geo:word-law} to words in their $g$ functions.  The error
cost after multiplication by $h$ is $O(R\eta)$, since all ranges lie
in a common space of dimension $O(R/h)$.  The difference of two such
self-adjoint functions need not commute with either factor.
Approximate the continuous function $|x|$ on its bounded spectral
interval by polynomials zero at zero, using that same finite range.
This proves, for any real $c,c_0$,
\[
\begin{gathered}
 h\Tr\left|c\,g(\mathcal A_h^{d,R})
                   -c_0g(\mathcal A_{0,h}^{d,R})\right|
\\
 \longrightarrow\frac1{2\pi}\int
 \operatorname{tr}|c\,g(\chi B)-c_0g(\chi B_0)|\,\dd\xi\dd s.
\end{gathered}
\]
On removing cutoffs, the low-frequency symbol cost is $O(d)$ because
$g$ is bounded.  Since $g(b)=O(b^2)$ near zero and
$\norm{B}\leq w_{\max}/(2|\xi|)$, the high-frequency symbol has an
integrable $C|\xi|^{-2}$ bound, uniformly in $\chi$.
Use \eqref{sgt:geo:cutoff-removal} on the operator side and
$|\Tr|X|-\Tr|Y||\leq\norm{X-Y}_1$ for absolute traces.
Finally \eqref{sgt:geo:actual-model} transfers all conclusions to the
actual geometry.  In particular
\begin{gather}
 h\Tr[(M_q\otimes V)g(X_h)]
\\
 \longrightarrow\frac1{2\pi}\int
              q(s)\operatorname{tr}[Vg(B)]\,\dd\xi\dd s,
 \label{sgt:geo:general-trace}\\
 h\Tr|c\,g(X_h)-c_0g(X_{0,h})|
\\
 \longrightarrow\frac1{2\pi}\int
             \operatorname{tr}|c\,g(B)-c_0g(B_0)|\,\dd\xi\dd s,
 \label{sgt:geo:general-absolute}
\end{gather}
for smooth $q$ and fixed matrix $V$, with the analogous reference law.
Fourier insertions described above are also allowed.  All estimates
include the singular low-frequency end and the infinite high-frequency
tail.  For heat a stronger symbol bound is
\begin{equation}
 0\leq F(s,\xi)\leq
       \ee^{-\tau(2|\xi|/w_{\max})^\gamma}I_M.
 \label{sgt:geo:heat-majorant}
\end{equation}
This is spectral calculus against a scalar multiple of $I_M$, not
operator monotonicity of the exponential.

For later receiver constructions it is useful to include sharp
Fourier bands explicitly.  Let $\Pi_{h;d,R}$ project onto
$d<|hk_n|<R$ in every component, where $0\leq d<R<\infty$.
For either heat operator and any fixed matrix $V$,
\begin{equation}
 h\Tr[\Pi_{h;d,R}V A_h(\tau)]
 \longrightarrow\frac1{2\pi}\int_\Gamma
            \int_{d<|\xi|<R}\operatorname{tr}(VF)\,\dd\xi\dd s.
 \label{sgt:geo:band-law}
\end{equation}
To see this, approximate the band indicator by a smooth Fourier
multiplier differing only in intervals of total length $O(\delta)$
around its endpoints.  The difference has norm at most one and rank
$O(\delta/h)+O(1)$.  Because $\norm{A_h}\leq1$, its trace against
$VA_h$ costs at most $\norm{V}$ times this rank; after multiplying by
$h$ this is $O(\delta)+O(h)$.  The symbol edge cost has the same bound.
Let $h\downarrow0$ and then $\delta\downarrow0$.  The possible
$n=0$ ambiguity costs only $M$ dimensions.  The smooth insertion
argument applies to the full heat by \eqref{sgt:geo:cutoff-removal}
and \eqref{sgt:geo:actual-model}.  Subtracting the $|hk_n|\leq R$ law
from the complete trace law proves the $|hk_n|>R$ law as well.
In particular its $V=Q$ or $V=I_M$ heat coefficient tends to zero as
$R\uparrow\infty$ by \eqref{sgt:geo:heat-majorant}.

Taking $g=f_{\tau,\gamma}$ now proves \eqref{sgt:geo:heat-law} and
the numerator in \eqref{sgt:geo:q-law}; the reference integral is
elementary and uses $\int w\dd s=1$.  If $c_h=hZ_h$ and $c_{0,h}=hZ_{0,h}$,
then
$\rho_h-\rho_{0,h}=h(A_h/c_h-A_{0,h}/c_{0,h})$.
Replacing these denominators by their positive limits costs $o(1)$
in trace norm, since each heat operator is positive.  Apply
\eqref{sgt:geo:general-absolute} with the limiting denominators to
obtain \eqref{sgt:geo:distance}.
For $M\geq2$, $F$ is positive definite at every $\xi\ne0$, so
$\operatorname{tr}(QF)>0$.  The variational characterization of trace
distance with the effect $Q$ gives $p_Q\leq D$, because
$Q\rho_{0,h}=0$.  Finally positive matrices $X,Y$ satisfy
$\operatorname{tr}|X-Y|\leq\operatorname{tr}X+\operatorname{tr}Y$,
with equality exactly when their supports are orthogonal.  This follows
by evaluating $X-Y$ on its positive and negative spectral projections:
equality forces each positive matrix into the corresponding disjoint
subspace.  Here $F/C_\gamma$ is positive definite and $F_0/C_{\rm col}$
is nonzero, so the inequality is strict pointwise.  The integral of
their summed traces is two, proving $D<1$.

\section{Growing-index counting, a strict deficit, and transverse rotation}
\label{sgt:geo:counting-section}
\label{sgt:geo:structure}

For fixed $s$, order the eigenvalues of $B(s,q)$ as
$b_1>\cdots>b_M>0$ and put $e_j=1/b_j$, with simplicity justified
below.  The free one-dimensional resolvent gives
\begin{equation}
 x^*\frac{T(q)}{2q}x
 =\frac1{2\pi}\int_{\mathbb R}
 \frac{|\sum_j\sqrt{\alpha_j}x_j\ee^{ita_j}|^2}{t^2+q^2}\,\dd t.
 \label{sgt:geo:free-resolvent}
\end{equation}
Subtracting this formula at $q_2>q_1>0$ proves a positive definite
difference.  Thus all ordered $b_j(s,q)$ strictly decrease and all
$e_j(s,q)$ strictly increase.  This is a statement about $T(q)/(2q)$;
the eigenvalues of $T(q)$ itself cannot all decrease since its trace
is one.

Set $\Delta_i=a_{i+1}-a_i>0$ and $A_\alpha=\operatorname{diag}\alpha_i$.
The exact energy matrix is
\begin{equation}
 B(s,q)^{-1}=w(s)^{-1}A_\alpha^{-1/2}J(q)A_\alpha^{-1/2},
 \label{sgt:geo:path-inverse}
\end{equation}
where, for $M\geq2$,
\begin{gather*}
 J_{11}=q[1+\coth(q\Delta_1)],\\
 J_{MM}=q[1+\coth(q\Delta_{M-1})],\\
 J_{ii}=q[\coth(q\Delta_{i-1})+\coth(q\Delta_i)]\\
(1<i<M),\\
 J_{i,i+1}=J_{i+1,i}=-q\operatorname{csch}(q\Delta_i),
\end{gather*}
and the other entries vanish.  For $M=1$ use $J(q)=2q$.
An energy derivation fixes every constant.  Minimize
$\int_{\mathbb R}(|u'|^2+q^2|u|^2)\dd a$ subject to $u(a_i)=x_i$.
On the interval of length $\Delta_i$, the hyperbolic interpolation
solving $-u''+q^2u=0$ has energy
\[
 q\coth(q\Delta_i)(|x_i|^2+|x_{i+1}|^2)
 -2q\operatorname{csch}(q\Delta_i)
                 \operatorname{Re}(\overline{x_i}x_{i+1}).
\]
The two exterior half-lines contribute $q|x_1|^2+q|x_M|^2$.
The reproducing kernel for the same energy is
$\ee^{-q|a-b|}/(2q)$, so its evaluation Gram matrix has inverse $J(q)$.
Conjugating by the masses proves \eqref{sgt:geo:path-inverse}.

As $q\downarrow0$, $J(q)$ tends to the nearest-neighbor path Laplacian
$J(0)$ with quadratic form
$\sum_i|x_{i+1}-x_i|^2/\Delta_i$.  Write
\begin{equation}
 \Lambda_\alpha=A_\alpha^{-1/2}J(0)A_\alpha^{-1/2},
 \qquad 0=\ell_1<\ell_2<\cdots<\ell_M.
 \label{sgt:geo:thresholds}
\end{equation}
The kernel is spanned by $v$.  Positivity of the other eigenvalues
follows from the quadratic form, and simplicity follows from the
irreducible symmetric tridiagonal recurrence: an eigenvector is
determined by its first component, which cannot vanish.  The same
argument applies for $q>0$ and proves simplicity in the definition of
$e_j$.  Consequently
\[
\begin{gathered}
 e_j(s,0+)=\ell_j/w(s),\\
 e_j(s,q)\geq2q/w(s),
 \\
 e_j(s,q)\longrightarrow\infty\\
(q\to\infty).
\end{gathered}
\]
The bound uses $b_j\leq\operatorname{tr}B=w/(2q)$.
Every branch therefore has exactly one positive-width energy cutoff
at energy $E$ when $w(s)E>\ell_j$, and none when $w(s)E\leq\ell_j$.

For $g>0$ let $-\kappa_j(g)^2$ be the negative eigenvalues of the
attractive delta Hamiltonian
\[
 -\frac{\dd^2}{\dd a^2}-g\sum_i\alpha_i\delta_{a_i}
 \quad\hbox{on }L^2(\mathbb R),
\]
defined by the form
$\int|u'|^2\dd a-g\sum_i\alpha_i|u(a_i)|^2$ on $H^1(\mathbb R)$.
Point evaluation is infinitesimally form bounded: for instance
$|u(a)|^2\leq\norm{u}_2\norm{u'}_2$, followed by Young's inequality,
gives an arbitrarily small coefficient of $\norm{u'}_2^2$.
Thus the form is closed and lower bounded.  Its negative index is at
most $M$, because the perturbation factors through $M$ evaluations.
Factoring those evaluations through
$(-\dd^2/\dd a^2+q^2)^{-1/2}$ shows that the negative index after adding
$q^2$ equals the number of eigenvalues of $gT(q)/(2q)$ above one;
the nonzero spectra of $XX^*$ and $X^*X$ coincide.  Equality with one
also gives the eigenspace at energy $-q^2$ by the weak equation.
The top eigenvalue diverges as $q\downarrow0$, so a lowest negative
state always exists.  We set $\kappa_j(g)=0$ for absent branches.
The preceding strictly monotone energy branches imply
\begin{equation}
 \#\{j:b_j(s,q)\geq E^{-1}\}
 =\#\{j:\kappa_j(w(s)E)\geq q\}\quad(q>0),
 \label{sgt:geo:birman-schwinger}
\end{equation}
including the equality threshold.  In particular the number of
negative states is $\#\{j:\ell_j<g\}$; equality at a threshold does
not add a negative state.

The actual growing-index count is
\begin{equation}
\begin{gathered}
 h\#\{k:\lambda_k(L_h)\leq E/h\}
 \longrightarrow
\\
 \mathcal N_\alpha(E)=\frac1\pi\int_\Gamma
                    \sum_{j=1}^M\kappa_j(w(s)E)\,\dd s,
 \\
 E>0.
 \label{sgt:geo:counting}
\end{gathered}
\end{equation}
To pass rigorously from continuous tests to this cutoff, approximate
$1_{[1/E,\infty]}$ on $[0,\infty]$ from above and below by continuous
functions zero near zero.  They belong to the class used in
\eqref{sgt:geo:general-trace}.  For each $s$, strict monotonicity gives
at most one positive $q$ with $b_j(s,q)=1/E$.  Fubini proves that the
phase-space level has measure zero.  The transition interval is
supported in a bounded frequency region, because
$b_j\leq w_{\max}/(2q)$.  Dominated convergence squeezes the two
phase-space approximations together.  Integrating
\eqref{sgt:geo:birman-schwinger} then gives \eqref{sgt:geo:counting}.
This proof controls the whole growing count, rather than a sum of
fixed-eigenvalue approximations.  For $L_h^\gamma$ below $U/h^\gamma$,
replace $E$ by $U^{1/\gamma}$.

\subsection{A self-contained finite-profile form of the classical sharp bound}
\label{sgt:geo:hlt}

The inequality $\sum_j\kappa_j(g)\leq g/2$, with equality only for a
single point mass, is the sharp measure Lieb--Thirring result of
Hundertmark, Lieb, and Thomas, Theorem~5 of
Ref.~\cite{HundertmarkLiebThomas1998}.
Here is a finite-matrix derivation sufficient for the present profile,
so the counting application does not require a measure approximation.

Write $\mathcal L(q)=gT(q)$ and
$\Sigma_j(q)=\sum_{i=1}^j\lambda_i(\mathcal L(q))$.
For $q_2>q_1\geq0$, set $t=q_2-q_1$ and
$U_\zeta=\operatorname{diag}(\ee^{i\zeta a_i})$.  The Fourier transform
of the Cauchy density gives the exact unitary average
\[
 \mathcal L(q_2)=\int_{\mathbb R}U_\zeta\mathcal L(q_1)U_\zeta^*
                   \frac{t\,\dd\zeta}{\pi(t^2+\zeta^2)}.
\]
The sum of the $j$ largest eigenvalues of a Hermitian matrix is
$\sup\{\operatorname{tr}(EX):0\leq E\leq I,\operatorname{tr}E=j\}$,
which is convex and unitarily invariant.  Therefore
$\Sigma_j(q_2)\leq\Sigma_j(q_1)$.
For $M\geq2$ the top eigenvalue is strictly decreasing: the positive
matrix $\mathcal L(q_2)$ has a strictly positive top eigenvector, and
every off-diagonal entry strictly increases on changing $q_2$ to
$q_1$.  Its Rayleigh quotient strictly increases as well.

Order $\kappa_1>\kappa_2>\cdots$ over the existing states and append
zeros to length $M$.  The energy simplicity established above gives
the strict first inequality.  The Birman--Schwinger equation gives
$2\kappa_j=\lambda_j(\mathcal L(\kappa_j))$ for existing states.
For absent states this is also true: $\mathcal L(0)=gP$ has all
eigenvalues after the first equal to zero.  Telescoping partial sums
gives the exact identity
\begin{equation}
\begin{gathered}
 2\sum_{j=1}^M\kappa_j
\\
 =\Sigma_M(\kappa_M)
    +\sum_{j=1}^{M-1}
             [\Sigma_j(\kappa_j)-\Sigma_j(\kappa_{j+1})]
\\
 \leq g+\Sigma_1(\kappa_1)-\Sigma_1(\kappa_2)<g.
 \label{sgt:geo:strict-local}
\end{gathered}
\end{equation}
The trace $\Sigma_M(q)$ equals $g$ at every $q$.  Every bracket is
nonpositive, and the first is strictly negative, including when only
one negative state exists and $\kappa_2=0$.  For $M=1$ directly
$\kappa_1=g/2$.  Repeated heights, if initially allowed, are first
merged; equality is therefore exactly the single-point case.
The sharp inequality and its rigidity are classical; their use here
is after the actual physical-symbol identification.

Define $\delta_\alpha(g)=g/2-\sum_j\kappa_j(g)$.  Equations
\eqref{sgt:geo:counting} and \eqref{sgt:geo:strict-local} imply
\begin{equation}
\begin{gathered}
 0<\mathcal N_\alpha(E)<\frac E{2\pi},\\
 \frac E{2\pi}-\mathcal N_\alpha(E)
       =\frac1\pi\int_\Gamma\delta_\alpha(w(s)E)\,\dd s
 \\
(M\geq2).
 \label{sgt:geo:count-deficit}
\end{gathered}
\end{equation}
The symbol heat is the Stieltjes integral against this counting
function.  Integration by parts, or positive layer-cake integration,
gives
\begin{gather}
 C_\gamma(\tau)=\tau\gamma\int_0^\infty
             E^{\gamma-1}\ee^{-\tau E^\gamma}
             \mathcal N_\alpha(E)\,\dd E,\notag\\
 C_{\mathrm{col},\gamma}(\tau)-C_\gamma(\tau)
\\
 =\frac{\tau\gamma}{\pi}\int_\Gamma\int_0^\infty
 E^{\gamma-1}\ee^{-\tau E^\gamma}\delta_\alpha(w(s)E)
                   \,\dd E\dd s>0.
 \label{sgt:geo:partition-deficit}
\end{gather}
There are no boundary terms because
$0\leq\mathcal N_\alpha(E)\leq E/(2\pi)$ and
$\mathcal N_\alpha(0)=0$.  This proves strict partition deficit for
every finite $\tau>0$ and every $\gamma>0$, with equality only for the
merged single-point profile.

\subsection{Algebraic transverse occupation from the lowest branch}
\label{sgt:geo:rotation}

Define profile and density statistics
\begin{equation}
\begin{gathered}
 d_i=\sum_j\alpha_j|a_i-a_j|,\\
 G=\sum_i\alpha_i d_i,\\
 V=\sum_i\alpha_i(d_i-G)^2,\\
 I_k=\int_\Gamma w(s)^k\,\dd s.
 \label{sgt:geo:profile-statistics}
\end{gathered}
\end{equation}
Let $\mathsf D_{ij}=\sqrt{\alpha_i\alpha_j}|a_i-a_j|$.
The finite matrix $T(q)=P-q\mathsf D+O(q^2)$ is analytic at zero.
Its simple top eigenpair, normalized by $v^*u(q)>0$, therefore obeys
\begin{equation}
\begin{gathered}
 \mu(q)=1-Gq+O(q^2),\\
 u(q)=v-qQ\mathsf Dv+O(q^2),\\
 \norm{Qu(q)}^2=Vq^2+O(q^3).
 \label{sgt:geo:eigenvector-rotation}
\end{gathered}
\end{equation}
These coefficients also follow directly by differentiating
$T(q)u(q)=\mu(q)u(q)$: the $v$ component gives $\mu'(0)=-G$;
the $Q$ component and normalization give $u'(0)=-Q\mathsf Dv$;
and $(Q\mathsf Dv)_i=\sqrt{\alpha_i}(d_i-G)$.
For small $g$ only the first bound state exists, by
\eqref{sgt:geo:thresholds}.  Its equation $2\kappa=g\mu(\kappa)$
and the implicit function theorem give
\begin{equation}
 \kappa(g)=\frac g2-\frac G4g^2+O(g^3),\qquad
 \kappa'(g)=\frac12-\frac G2g+O(g^2).
 \label{sgt:geo:weak-coupling}
\end{equation}
For $M=1$ the same expressions hold with $G=V=0$.
The mean-distance coefficient $G$ is the familiar one-dimensional
weak-coupling statistic \cite{Simon1976,Simon2022}; see
Theorem~8.5 and Eq.~(9.24) of Ref.~\cite{Simon2022}.
The finite atomic coefficient used here has been derived directly;
its occurrence alone is not a new scalar invariant.

Choose a small fixed energy $E_0$ below $\ell_2/w_{\max}$ when $M\geq2$,
and small enough for \eqref{sgt:geo:weak-coupling} uniformly over $s$.
On the lowest branch put $q=\kappa(w(s)E)$.  Then
\begin{align}
 \frac{\dd q}{\dd E}
 &=\frac{w(s)}2-\frac{Gw(s)^2}2E+O(E^2),\notag\\
 \norm{Qu(\kappa(w(s)E))}^2\frac{\dd q}{\dd E}
 &=\frac{Vw(s)^3}{8}E^2+O(E^3),
 \label{sgt:geo:jacobian}
\end{align}
with uniform errors.  The Taylor formulas are used only on that fixed
small interval.  To bound the rest of every branch, put
$a_*=\ell_2/w_{\max}>0$.  For $j\geq2$,
$e_j(s,q)\geq a_*$ and $e_j(s,q)\geq2q/w_{\max}$, whence
\begin{equation}
\begin{gathered}
 C_{\mathrm{exc},\gamma}(\tau):=
 \frac1\pi\int_\Gamma\int_0^\infty
       \sum_{j=2}^M\ee^{-\tau e_j(s,q)^\gamma}\,\dd q\dd s
\\
 \leq C\tau^{-1/\gamma}\ee^{-\tau a_*^\gamma/2}.
 \label{sgt:geo:excited-tail}
\end{gathered}
\end{equation}
This follows from the integrable product bound
$\ee^{-\tau e_j^\gamma}
\leq\ee^{-\tau a_*^\gamma/2}
\ee^{-(\tau/2)(2q/w_{\max})^\gamma}$.
The portion of the lowest branch with $E\geq E_0$ has the identical
bound with $a_*$ replaced by $E_0$, by integrating in $q$.
Thus the infinite tails are exponentially small, including their
phase volume.

Now use
$\int_0^\infty E^r\ee^{-\tau E^\gamma}\dd E
=\Gamma((r+1)/\gamma)\tau^{-(r+1)/\gamma}/\gamma$ for $r>-1$.
Integrating \eqref{sgt:geo:jacobian} and its uniformly bounded
remainders gives
\begin{gather}
 C_\gamma(\tau)
 =\frac{G_\gamma}{2\pi}\tau^{-1/\gamma}
 -\frac{GI_2\Gamma(1+2/\gamma)}{4\pi}\tau^{-2/\gamma}
 +O(\tau^{-3/\gamma}),
 \label{sgt:geo:cold-partition}\\
 C_{Q,\gamma}(\tau)
 =\frac{VI_3\Gamma(1+3/\gamma)}{24\pi}\tau^{-3/\gamma}
 +O(\tau^{-4/\gamma}),\notag\\
 p_{Q,\gamma}(\tau)
 =\frac{VI_3\Gamma(1+3/\gamma)}{12G_\gamma}
                       \tau^{-2/\gamma}
 +O(\tau^{-3/\gamma}).
 \label{sgt:geo:cold-q}
\end{gather}
For $\gamma=1$ the latter leading coefficient is $VI_3/2$;
for $\gamma=1/2$ it is $30VI_3$.  The excited fraction in
\eqref{sgt:geo:excited-tail}, divided by $C_\gamma$, is exponentially
small.  The algebraic term in \eqref{sgt:geo:cold-q} comes instead
from the rotation of the lowest symbol eigenspace relative to the
fixed comparison $P$.  These are limiting symbol branches, not a
claim of exact global band labels for individual curved eigenstates.

The vanishing cases of $V$ can be determined exactly.  With
$F_i=\sum_{j\leq i}\alpha_j$,
\[
 d_{i+1}-d_i=\Delta_i(2F_i-1).
\]
Since all weights and gaps are positive, $V=0$ means all $d_i$ coincide.
For $M\geq2$ this forces every $F_i$ to equal $1/2$, which is possible
only for $M=2$ with equal weights.  For that profile $T(q)$ has constant
row sums and $v$ is its top eigenvector at every $q$; the lowest branch
has exactly zero $Q$ weight and $C_Q=C_{\rm exc}$ is exponentially
small.  For $M=1$, $Q=0$.  Every unequal two-layer profile and every
profile with $M\geq3$ has $V>0$ and the algebraic leading term is
strictly positive.

\section{One-excitation scalar agreement and recoverable information}
\label{sgt:geo:information-section}

For each fixed integer $p\geq0$ the test
$g_{p,\tau,\gamma}(b)=(\tau b^{-\gamma})^p\ee^{-\tau b^{-\gamma}}$
with value zero at $b=0$ belongs to the class already proved.
It is exponentially flat at zero and has a finite limit at infinity,
equal to one for $p=0$ and zero for $p>0$.
Thus, independently of differentiating any asymptotic remainder,
\begin{equation}
\begin{gathered}
 h\Tr[(\tau X_h^{-\gamma})^pA_h(\tau)]
\\
 \longrightarrow C_{p,\gamma}(\tau):=\\
 \frac1{2\pi}\int_\Gamma\int_{\mathbb R}
\\
 \operatorname{tr}[(\tau B^{-\gamma})^p\ee^{-\tau B^{-\gamma}}]
                         \,\dd\xi\dd s,
 \\
 C_{0,\gamma}=C_\gamma.
 \label{sgt:geo:weighted-heat}
\end{gathered}
\end{equation}
Products with inverse powers here denote their bounded spectral
functions, zero on the inverse kernel.  The same result holds for
each fixed real $p\geq0$ if desired.  The high-frequency estimate is
$x^p\ee^{-x}\leq C_p\ee^{-x/2}$; the low-frequency estimate is simply
boundedness of this scalar test.  These also justify differentiation
of the limiting symbol integrals:
$C_{1,\gamma}=-\tau C_\gamma'$ and
$C_{2,\gamma}=\tau^2C_\gamma''$.

For every fixed $r>0$, $r\ne1$, the exact Gibbs identity
$\Tr\rho_h^r=Z_h(r\tau)/Z_h(\tau)^r$ gives the Renyi entropy
\begin{equation}
 S_r(\rho_h)=\log(1/h)+
 \frac{\log C_\gamma(r\tau)-r\log C_\gamma(\tau)}{1-r}+o(1).
 \label{sgt:geo:renyi}
\end{equation}
The von Neumann entropy is treated separately, using its exact Gibbs
identity and the $p=1$ case of \eqref{sgt:geo:weighted-heat}:
\begin{equation}
 S(\rho_h)=\log(1/h)+\log C_\gamma(\tau)
                 +\frac{C_{1,\gamma}(\tau)}{C_\gamma(\tau)}+o(1).
 \label{sgt:geo:entropy}
\end{equation}
Their finiteness at fixed $h$ follows from the spectral exponential
estimate given after \eqref{sgt:geo:unitary}; no dimension-free entropy
continuity is used.  Likewise
\begin{equation}
 \beta\langle L_h^\gamma\rangle\longrightarrow
       \frac{C_{1,\gamma}}{C_\gamma},\qquad
 \beta^2\operatorname{Var}(L_h^\gamma)\longrightarrow
       \frac{C_{2,\gamma}}{C_\gamma}
          -\left(\frac{C_{1,\gamma}}{C_\gamma}\right)^2.
 \label{sgt:geo:energy}
\end{equation}
All reference formulas replace $C_\gamma$ by $C_{\mathrm{col},\gamma}$.
No uniformity in entropy order approaching zero or one, or in an
unbounded family of moments, is asserted.

To take either endpoint in $\tau$, set
$\xi=w(s)u/\tau^{1/\gamma}$.  The heat matrix becomes
\[
 f_{(2|u|)^\gamma,\gamma}
       \bigl(T(w(s)|u|/\tau^{1/\gamma})\bigr).
\]
Its eigenvalues are bounded by $\ee^{-(2|u|)^\gamma}$ since
$\norm{T(q)}\leq1$.  This is an integrable common majorant independent
of $s$ and $\tau$; weighted energy tests have the analogous majorant
$C_p\ee^{-(2|u|)^\gamma/2}$.  For $u\ne0$, distinct heights give
$T(w|u|/\tau^{1/\gamma})\to\operatorname{diag}\alpha_i$ as
$\tau\downarrow0$.  As $\tau\uparrow\infty$ the same matrix tends to
$P$, and continuity of the spectral function at zero includes its
kernel.  Dominated convergence applies to traces and matrix trace
norms, with the extra spatial weight $w(s)$ integrating to one.
Writing
$A_{\alpha,\gamma}(u)=\operatorname{diag}
(\ee^{-(2|u|/\alpha_i)^\gamma})$ and
$b_\gamma(u)=\ee^{-(2|u|)^\gamma}$, this proves
\begin{gather}
 \frac{C_\gamma(\tau)}{C_{\mathrm{col},\gamma}(\tau)}\longrightarrow1,
 \\
p_{Q,\gamma}(\tau)\longrightarrow1-\sum_i\alpha_i^2,
 \notag\\
 D_\gamma(\tau)\longrightarrow
 D^*_{\alpha,\gamma}:=\frac1{G_\gamma}\int_0^\infty
       \operatorname{tr}|A_{\alpha,\gamma}(u)-b_\gamma(u)P|\,\dd u
 \\
(\tau\downarrow0),
 \label{sgt:geo:hot-endpoint}\\
 \frac{C_\gamma(\tau)}{C_{\mathrm{col},\gamma}(\tau)}\longrightarrow1,
 \\
p_{Q,\gamma}(\tau)\longrightarrow0,\\
 D_\gamma(\tau)\longrightarrow0\\
(\tau\uparrow\infty).
 \label{sgt:geo:cold-endpoint}
\end{gather}
For normalization, the integral over $\mathbb R$ of the $i$th
diagonal high-resolution heat is $\alpha_iG_\gamma$, while the
reference integral is $G_\gamma$.  Thus the normalized trace-distance
prefactor is $1/(2G_\gamma)$ on the full line, or $1/G_\gamma$ on the
positive half-line, as written.  Positive limiting denominators and
the common majorant justify their replacement inside the integral.

The same substitution with $x=(2u/\alpha_i)^\gamma$ gives
\begin{equation}
 \frac{C_{p,\gamma}(\tau)}{C_\gamma(\tau)}
 \longrightarrow\frac{\Gamma(p+1/\gamma)}{\Gamma(1/\gamma)}
 \quad(\tau\downarrow0).
 \label{sgt:geo:moment-endpoint}
\end{equation}
Indeed the positive-half-line integral of $x^p\ee^{-x}$ is
$\alpha_i\Gamma(p+1/\gamma)/(2\gamma)$, and the weights sum to one.
The dimensionless mean and variance in \eqref{sgt:geo:energy} both
tend to $1/\gamma$.  From \eqref{sgt:geo:renyi}, the partition ratio
at $\tau$ and at fixed $r\tau$, and \eqref{sgt:geo:entropy}, one obtains
\begin{equation}
\begin{gathered}
 \lim_{\tau\downarrow0}\lim_{h\downarrow0}
     [S_r(\rho_h)-S_r(\rho_{0,h})]=0,
 \\
 \lim_{\tau\downarrow0}\lim_{h\downarrow0}
     [S(\rho_h)-S(\rho_{0,h})]=0.
 \label{sgt:geo:scalar-agreement}
\end{gathered}
\end{equation}
These are additive entropy differences.  The limits do not merely
compare two divergent entropies by their ratio.

For equal weights $\alpha_i=1/M$, put
$a(u)=\ee^{-(2M u)^\gamma}$ and $b(u)=\ee^{-(2u)^\gamma}$ on $u>0$.
Since $a\leq b$, the eigenvalues of $aI_M-bP$ give
$\operatorname{tr}|aI_M-bP|=b+(M-2)a$.
The integrals of $b$ and $a$ are respectively $G_\gamma/2$ and
$G_\gamma/(2M)$.  Hence
\begin{equation}
 \lim_{\tau\downarrow0}\lim_{h\downarrow0}
       \frac12\norm{\rho_h-\rho_{0,h}}_1
 =\lim_{\tau\downarrow0}\lim_{h\downarrow0}\Tr(Q\rho_h)
 =1-\frac1M.
 \label{sgt:geo:equal-information}
\end{equation}
The binary decision ``multilayer on a $Q$ click, collapsed otherwise''
has equal-prior error $(1-\Tr(Q\rho_h))/2$, since the reference never
clicks.  Its iterated limit is $1/(2M)$, equal to the single-copy
Helstrom optimum from \eqref{sgt:geo:equal-information}.
For any fixed finite number $N$ of independent copies its any-click
error tends to $1/(2M^N)$; collective optimality is not needed for this
last statement.  For unequal weights the $Q$ limit
$1-\sum_i\alpha_i^2$ remains a valid lower bound on the distance,
but need not equal the full matrix absolute-value integral.

For completeness that unequal-weight integral has a scalar
characterization.  The rank-one subtraction
$A_{\alpha,\gamma}(u)-b_\gamma(u)P$ has exactly one negative eigenvalue
$-r(u)$ for $u>0$, where $r(u)>0$ is the unique solution of
\[
 b_\gamma(u)\sum_i
 \frac{\alpha_i}{\ee^{-(2u/\alpha_i)^\gamma}+r(u)}=1.
\]
The left side at $r=0$ is strictly larger than one for $M\geq2$;
it decreases continuously to zero.  The rank-one eigenvector equation
gives the displayed condition and shows that there are no other
negative eigenvalues.  The signed trace integrates to zero, so
$D^*_{\alpha,\gamma}=(2/G_\gamma)\int_0^\infty r(u)\dd u$.
In particular $1-\sum_i\alpha_i^2\leq D^*_{\alpha,\gamma}<1$,
the strict upper bound following from positive definiteness and the
strict trace-norm inequality used above.

Every finite-accuracy consequence has the same quantifiers.  First
choose a finite profile, an accuracy, and a finite list of entropy
orders, energy moments, and effects.  Choose a sufficiently small
$\tau$ for their limiting functions, and then a sufficiently small
$h$ at that fixed $\tau$.  The displayed agreements and distinctions
then hold simultaneously to that accuracy.  A large but fixed $M$
can be selected before these limits.  This gives no numerical onset
or uniform prescription for $M(h)$ or $\tau(h)$.

Finally, a physical bounded Lipschitz multiplication detector $q$ on
the tube pulls back to
$\operatorname{diag}(q(T_{i,h}(s)))$.  It differs in norm from
$q(\mathbf r(s))I_M$ by at most
$h\max_i|a_i|\operatorname{Lip}(q)$.
Smooth spatial insertions in \eqref{sgt:geo:general-trace}, followed
by the same endpoint substitution, give the common limiting
tangential density $w(s)\dd s$ at both endpoints; uniform approximation
extends this to continuous bounded spatial tests on the curve.
The large $Q$ distinction requires coherent comparison of sheet
amplitudes in the declared alignment.  Neither that unitary
identification nor the one-excitation entropy calculation supplies
a common laboratory vacuum, an equilibration mechanism, or a
many-oscillator thermodynamic assertion.

% Included by the standalone SG-T manuscript; no independent preamble.
\section{Controlled preparation of the actual one-excitation states}
\label{sgt:prep:appendix}

Throughout this section the curve, its interior tube, the positive smooth
density $w$, the distinct layer heights $a_i$, their positive weights
$\alpha_i$, and the finite layer number $M$ are fixed. We normalize
$\int_\Gamma w\,\dd s=\sum_i\alpha_i=1$ and use
\[
\begin{gathered}
 \mathcal H=L^2(\Gamma,\dd s)\otimes\mathbb C^M,\\
 v=(\sqrt{\alpha_1},\ldots,\sqrt{\alpha_M})^T,\\
 P=vv^*,\\
 Q=I-P.
\end{gathered}
\]
Here and below $P,Q$ also denote their tensor extensions to $\mathcal H$.
Let $R_\varepsilon$ be the actual positive Green inverse after conjugation
from the pushforward masses $\alpha_iw(s)\dd s$ on
$T_{i,\varepsilon}(s)=\gamma(s)+\varepsilon a_i n(s)$.
The reference is $R_0=R_{\rm base}\otimes P$.
The physical eigenvalues of $L_h=R_h^{-1}$, $h\in\{0,\varepsilon\}$,
are $\lambda_j(L_h)=r_j(R_h)^{-1}$, where $r_j$ lists the positive
eigenvalues in decreasing order. The reference alignment complement is
assigned heat zero; it is never inserted as a set of zero-energy modes.
Here $\varepsilon$ is the geometric separation denoted by $h$ in the
main text; the local index $h\in\{0,\varepsilon\}$ only selects the
reference or split candidate in this preparation section.
All estimates below are uniform only for sufficiently small
$\varepsilon>0$ with these fixed data.

\subsection{An actual-operator weak trace-ideal estimate}
\label{sgt:prep:weakS1}

We prove, including the cross-layer coupling, that
\begin{equation}
 r_j(R_h)\le \frac{K}{j},\qquad
 \lambda_j(L_h)\ge c j,\qquad c=K^{-1}>0.
 \label{sgt:prep:spectral-bound}
\end{equation}
Thus $R_h$ is uniformly in the weak trace ideal $\mathfrak S_{1,\infty}$.
The assertion for $h=0$ concerns only its nonzero physical spectrum.

\paragraph{Uniform diagonal regularity.}
Let $\ell$ be the length of $\Gamma$, parametrized by arclength.
On a fixed closed interval $0\le\varepsilon\le\varepsilon_*$, every
individual offset remains embedded, has speed bounded away from zero,
and has uniformly bounded derivatives. Its diagonal layer kernel is
\[
 R_{ii,\varepsilon}(s,t)=\alpha_i\sqrt{w(s)w(t)}
 G(T_{i,\varepsilon}(s),T_{i,\varepsilon}(t)).
\]
The regular part of the Dirichlet Green function is smooth on the tube
product, also with any fixed distant Dirichlet masks. The quotient
\[
 \frac{|T_{i,\varepsilon}(s)-T_{i,\varepsilon}(t)|}
 {(\ell/\pi)|\sin(\pi(s-t)/\ell)|}
\]
extends positively and smoothly across $s=t$ as a periodic function.
Indeed, near the diagonal division of the chord by $s-t$ gives the
integral of the nonvanishing tangent; division of the squared norms
followed by the positive square root proves smoothness. Away from the
diagonal compactness and uniform embeddedness give a positive lower
bound. All fixed orders of derivatives are uniformly bounded.
Subtracting $-(2\pi)^{-1}\log|2\sin(\pi(s-t)/\ell)|$ therefore leaves
a uniformly smooth family. For the nonzero Fourier frequencies
$k_n=2\pi n/\ell$, this logarithmic convolution has multiplier
$1/(2|k_n|)$. Its zero multiplier is bounded and immaterial to the
estimate. It maps $L^2$ to $H^1$ boundedly, as do the smooth remainders
and multiplication by $\sqrt w$. Consequently
\[
 \sup_{i,\varepsilon}\|R_{ii,\varepsilon}:L^2\longrightarrow H^1\|<\infty.
\]
The embedding $H^1(\Gamma)\hookrightarrow L^2(\Gamma)$ has singular
values bounded by $C/j$, by the same Fourier basis. Factoring each
diagonal block through this embedding yields
$s_j(R_{ii,\varepsilon})\le C_d/j$ with one common constant.
The base curve obeys the same estimate.

\paragraph{The cross blocks are controlled by positivity.}
For every positive operator $R$ on a sum of $M$ Hilbert spaces, let
$E_i$ inject the $i$th summand. Then
\[
\begin{gathered}
 \langle f,Rf\rangle
 =\left\|\sum_iR^{1/2}E_if_i\right\|^2
\\
 \le M\sum_i\|R^{1/2}E_if_i\|^2
 =M\sum_i\langle f_i,R_{ii}f_i\rangle.
\end{gathered}
\]
Hence, as a form inequality of bounded positive operators,
\begin{equation}
 0\le R_\varepsilon\le
 M\,\operatorname{diag}(R_{11,\varepsilon},\ldots,R_{MM,\varepsilon}).
 \label{sgt:prep:positive-block-bound}
\end{equation}
Each diagonal block is positive because it is a compression of
$R_\varepsilon$. The number of its eigenvalues exceeding $t>0$ is at
most $C_d/t$. The direct sum consequently has at most $MC_d/t$ such
eigenvalues. Min--max applied to
\eqref{sgt:prep:positive-block-bound} proves
\eqref{sgt:prep:spectral-bound}, after enlarging a fixed constant to
include the reference. No uniform smoothness of a cross-layer kernel
was assumed. Such a smoothness assertion would fail as the layers
collide. No constant uniform in growing $M$ or vanishing weights is
claimed.

\subsection{Heat tails and a common tensor encoding}
\label{sgt:prep:encoding}

Fix $\tau,\gamma>0$, set $\beta=\tau\varepsilon^\gamma$, and write
$A_h=\exp(-\beta L_h^\gamma)$ with the reference convention above.
Retain its first $J$ physical spectral terms in $A_{h,J}$, and put
$Z_h=\Tr A_h$, $Z_{h,J}=\Tr A_{h,J}$ and
$T_{h,J}=Z_h-Z_{h,J}$. For $a=\tau c^\gamma$, monotonicity and an
integral comparison give the explicit all-mode estimate
\begin{align}
 T_{h,J}
 &\le\sum_{j>J}\exp[-a(\varepsilon j)^\gamma]
 \le\frac1\varepsilon\int_{\varepsilon J}^\infty
                   \exp(-a x^\gamma)\,\dd x\notag\\
 &=\frac{\Gamma(1/\gamma,a(\varepsilon J)^\gamma)}
          {\varepsilon\gamma a^{1/\gamma}}.
 \label{sgt:prep:tail}
\end{align}
Here $\Gamma(s,y)=\int_y^\infty t^{s-1}e^{-t}\dd t$.
For fixed $s>0$ and $y\ge1$ it is at most $B_s e^{-y/2}$:
bound $t^{s-1}e^{-t/2}$ on $[1,\infty)$ and integrate the remaining
exponential. The actual joint heat law, Eq.~\eqref{sgt:geo:heat-law}, gives constants
$0<z_*\le z^*<\infty$ such that, eventually, under either hypothesis,
\begin{equation}
 \frac{z_*}{\varepsilon}\le Z_h\le\frac{z^*}{\varepsilon}.
 \label{sgt:prep:partition-scale}
\end{equation}
For any $0<\delta<1$, equations \eqref{sgt:prep:tail} and
\eqref{sgt:prep:partition-scale} allow
\begin{equation}
 J=\left\lceil\frac{\kappa_\delta}{\varepsilon}\right\rceil,
 \qquad
 \kappa_\delta=O\bigl([1+\log(1/\delta)]^{1/\gamma}\bigr),
 \qquad \frac{T_{h,J}}{Z_h}\le\delta.
 \label{sgt:prep:cutoff}
\end{equation}
The retained and discarded spectral supports are orthogonal, so
\begin{equation}
 \frac12\left\|\frac{A_h}{Z_h}
                  -\frac{A_{h,J}}{Z_{h,J}}\right\|_1
 =\frac{T_{h,J}}{Z_h}\le\delta.
 \label{sgt:prep:trace-distance}
\end{equation}
This is an asymptotic existence estimate with no certified finite onset.

Let $W\subset L^2(\Gamma,\dd s)$ be the span of all $M$ scalar
components of the $J$ retained target eigenvectors and the $J$ retained
base-curve eigenfunctions. Then $n=\dim W\le(M+1)J$ and
\begin{equation}
 \mathcal S=W\otimes\mathbb C^M,\qquad
 D=Mn\le M(M+1)J.
 \label{sgt:prep:dimension}
\end{equation}
Both retained state ranges lie in $\mathcal S$, which reduces every
fixed transverse matrix. For an orthonormal-basis isometry
$V_W:\mathbb C^n\to W$, set $V=V_W\otimes I_M$. Thus
\begin{equation}
\begin{gathered}
 \rho_{D,h}=\frac{V^*A_{h,J}V}{Z_{h,J}},\\
 Q_D=V^*QV=I_n\otimes Q,\\
 Q_D\rho_{D,0}=0.
 \label{sgt:prep:common-Q}
\end{gathered}
\end{equation}
The eigenbases of the two actual candidates need not agree.
For an original physical spatial multiplication effect $M_f$, its
encoded effect is precisely $V^*M_fV$, generally a dense matrix.
The physical spatial measurement class remains this compressed family;
arbitrary output-port effects do not acquire spatial locality by
being assigned a coordinate label. A fixed layer unitary sending $v$
to one output implements $Q_D$ in each of the $n$ longitudinal modes.
Combining one known $M$-component vector uses $M-1$ two-mode rotations,
hence $O(D)$ elements overall; the physical conversion into the modes
$W$, routing, and transmission have separate costs.

\paragraph{Weighted tails, rather than an entropy continuity shortcut.}
Let $y_j=\beta\lambda_j^\gamma$, $q=T_J/Z$, and
$m=Z^{-1}\sum_j y_je^{-y_j}$. With $m_J$ its retained normalized
version and $m_{\rm tail}=Z^{-1}\sum_{j>J}y_je^{-y_j}$,
\[
\begin{gathered}
 S(\rho)-S(\rho_J)=-\log(1-q)+m-m_J,
 \\
 |m-m_J|\le\frac{m_{\rm tail}+qm}{1-q}.
\end{gathered}
\]
For each fixed $p\ge0$, $y^pe^{-y}\le C_pe^{-y/2}$, so
\eqref{sgt:prep:tail} at $\tau/2$ controls each chosen moment tail.
The weighted joint law bounds the full moments. For any selected
positive R\'enyi order $r\ne1$, use also the cutoff estimate at
$r\tau$ and the identity
$S_r=(\log Z(r\beta)-r\log Z(\beta))/(1-r)$.
Thus a fixed finite list of thermal scalars has controlled truncation
error with the same dimension order, without an extra
$\log(1/\varepsilon)$ factor. Arbitrary preparation errors small only
in trace distance do not imply this stronger assertion at growing
dimension or for unbounded energy observables.

\subsection{Deliberate single-excitation preparation and its costs}
\label{sgt:prep:success}

On the common encoded space let $A_{D,h}=V^*A_{h,J}V$ and
$F_h=A_{D,h}^{1/2}$. Inject a heralded excitation into a uniformly
random input port, giving $I_D/D$, and apply the passive contraction
$F_h$. The surviving block, success probability, and conditional state
are exactly
\begin{equation}
 \sigma_h=\frac{A_{D,h}}D,\qquad
 s_h=\frac{Z_{h,J}}D,\qquad
 \frac{\sigma_h}{s_h}=\rho_{D,h}.
 \label{sgt:prep:filter-identity}
\end{equation}
In particular,
\begin{equation}
 s_h\ge
 \frac{(1-\delta)z_*}
 {M(M+1)(\kappa_\delta+\varepsilon)}.
 \label{sgt:prep:success-floor}
\end{equation}
At each fixed accuracy this is a positive lower bound independent of
small $\varepsilon$. It is preparation of an injected excitation,
not selection of one quantum from an equilibrium Bose sample. In that
different protocol the probability is $Z_1/Z_{\rm B}$ and can be
exponentially small on the joint scale.

The exact dilation
\[
 U_F=\begin{pmatrix}
 F&(I-F^2)^{1/2}\\ (I-F^2)^{1/2}&-F
 \end{pmatrix}
\]
is unitary because its blocks commute. Vacuum auxiliary ports realize
$F$ on the surviving one-excitation sector. Generic universal
compilation takes $O(D^2)$ two-mode elements and $O(D)$ depth
\cite{Clements2016}. Those counts do not establish constant
transmission: fixed loss per traversal at depth proportional to $D$
may make the collection probability exponentially small.

An alternative is to draw a retained eigenlabel with its normalized
Gibbs probability and prepare that known encoded vector. For
$c\in\mathbb C^D$, pad to $D'=2^d$, $d=\lceil\log_2D\rceil$.
At each node of a binary splitter tree, split power according to
the two descendant sums of $|c_i|^2$; at the leaves impose phases
$\arg c_i$. Induction gives the desired vector with $D'-1=O(D)$
splitters and depth $d$. Under the explicit assumption of a common
intensity transmission $\vartheta\in(0,1]$ per level, balanced paths
and calibrated phases, conditioning on survival preserves the vector
and
\[
 q_{\rm tree}\ge\vartheta^{\lceil\log_2D\rceil}
 \ge\vartheta D^{-\kappa},\qquad
 \kappa=-\frac{\log\vartheta}{\log2}.
\]
A mere lower bound on unequal path transmissions would not preserve
the vector. If known label-dependent transmissions $q_j$ preserve
each individual vector, sample label $j$ with probabilities
proportional to $p_j/q_j$. Its accepted mixture then has weights $p_j$
and acceptance $(\sum_jp_j/q_j)^{-1}\ge\min_jq_j$.
The eigenvectors, their certified approximation, $O(DJ)$ dense
coefficient storage, $O(D)$ settings per reconfiguration, settling and
reset are paid resources. This proves a polynomial transmission budget
only under the stated balanced-loss and polynomial-control assumptions.
It is neither an efficient continuum eigensolver nor an unknown-specimen
oracle: a programmed matrix is already known, whereas a specimen
experiment requires physical access to the unknown generator.

For $N$ accepted copies, source delivery and state-independent
collection factors $q_{\rm src},q_{\rm col}$ yield mean raw attempt
count $N/(s_hq_{\rm src}q_{\rm col})$. Failures, source reset and
filtering time remain charged. A final destructive click can certify
a completed measured trial; availability of a stored excitation
requires an actual nondestructive herald. A failed arrival is not a
measurement of $I-Q$. Acceptance rates, input labels, waiting times,
or resolved failure channels may themselves reveal the hypothesis;
claims about conditional normalized copies are not bounds on an
enlarged experiment retaining those records.

\subsection{Square-root Gibbs weights from ordinary loss}
\label{sgt:prep:subordination}

For $\gamma=1/2$ the classical subordination identity
\cite{Garrigos2015} is
\begin{equation}
\begin{gathered}
 e^{-\beta\sqrt L}=\int_0^\infty\ell_\beta(u)e^{-uL}\,\dd u,
 \\
 \ell_\beta(u)=\frac{\beta}{2\sqrt\pi\,u^{3/2}}
              e^{-\beta^2/(4u)}.
 \label{sgt:prep:subordination-identity}
\end{gathered}
\end{equation}
For completeness, set
$I(b)=\int_0^\infty e^{-x^2-b^2/x^2}\dd x$ for $b>0$.
Differentiation under the integral and the substitution $x=b/y$
give $I'(b)=-2I(b)$. Since $I(0)=\sqrt\pi/2$ by dominated
convergence, $I(b)=(\sqrt\pi/2)e^{-2b}$.
The substitution $x=\beta/(2\sqrt u)$ proves the scalar Laplace
identity in \eqref{sgt:prep:subordination-identity}; the spectral
theorem proves its operator form. Taking the scalar spectral value
zero proves that $\ell_\beta$ integrates to one.

Suppose a finite positive supported matrix $L$ has accessible linear
loss jumps $J_\mu=\sum_i b_{\mu i}a_i$ with $B^*B=L$.
Under the Lindblad convention
\[
 \dot\rho=\sum_\mu
 \left(J_\mu\rho J_\mu^*
       -\tfrac12\{J_\mu^*J_\mu,\rho\}\right),
\]
every jump from the one-excitation sector enters vacuum. Its surviving
block therefore evolves as
$\rho_{\rm surv}(u)=e^{-uL/2}\rho_{\rm surv}(0)e^{-uL/2}$,
the standard conditional-loss evolution \cite{PlenioKnight1998}.
Apply this with $L$ the restriction of the actual $L_h$ to its $J$
retained physical eigenvectors. Project the input $I_D/D$ onto that
support, counting its unused coordinates as initial failures.
For sampled duration $u$ the unnormalized survivor is $e^{-uL}/D$.
Pooling these unnormalized survivors with density $\ell_\beta$
gives $e^{-\beta\sqrt L}/D$ and hence exactly the state and success
in \eqref{sgt:prep:filter-identity}. Normalizing at each duration
before averaging would change the preparation. A commuting coherent
drift leaves the diagonal surviving block unchanged; a general
noncommuting drift needs its own control estimate.

\paragraph{A finite maximum duration.}
The sampled duration has infinite mean. Abort proposed draws $u>U$
before injection, or resample them classically. Their probability is
\begin{equation}
 q_U=\operatorname{erf}\!\left(\frac{\beta}{2\sqrt U}\right)
 \le\frac{\beta}{\sqrt{\pi U}}.
 \label{sgt:prep:duration-tail}
\end{equation}
Put $A=e^{-\beta\sqrt L}$,
$A_U=\int_0^U\ell_\beta(u)e^{-uL}\dd u$ and
$b=\Tr(A-A_U)$. Then $0\le A-A_U\le q_UI_J$ and
$b\le Jq_U\le Dq_U$. Adding and subtracting $A_U/\Tr A$ gives
\begin{equation}
 \frac12\left\|\frac A{\Tr A}-\frac{A_U}{\Tr A_U}\right\|_1
 \le\frac b{\Tr A}\le\frac{q_U}s,
 \qquad s=\frac{\Tr A}D.
 \label{sgt:prep:duration-bias}
\end{equation}
If $s_*\le s$ and additional state distance $\delta_p\in(0,1)$
is allowed, choose
\begin{equation}
 U\ge\frac{\beta^2}{\pi\delta_p^2s_*^2}.
 \label{sgt:prep:duration-cutoff}
\end{equation}
Counting all proposed draws gives success at least $(1-\delta_p)s$.
Resampling before injection gives per-injection success
$\Tr A_U/[D(1-q_U)]$ and mean classical draw count $1/(1-q_U)$.
Every injected filtering run lasts at most $U$. At
$\beta=\tau\sqrt\varepsilon$ and fixed accuracies the chosen $U$
is $O(\varepsilon)$ in the mathematical clock of $L$.

\paragraph{Exact finite mean with first-loss stopping.}
Alternatively, assume all losses are promptly and reliably monitored.
Sample the full duration distribution but stop at the first loss or
the sampled endpoint. In input eigenmode $j$, the loss time is
exponential with rate $\lambda_j$, independently of the sampled
duration. Tonelli's theorem and its Laplace transform give
\[
 \mathbb E\min(u,T_j)
 =\int_0^\infty\Pr(u>t)e^{-\lambda_jt}\dd t
 =\frac{1-e^{-\beta\sqrt{\lambda_j}}}{\lambda_j}.
\]
Including initial support failures with zero filtering time,
\begin{align}
 \mathbb ET_{\rm attempt}
 &=\frac1D\sum_{j\le J}
     \frac{1-e^{-\beta\sqrt{\lambda_j}}}{\lambda_j}
 \le\frac{\beta}{D}\sum_{j\le J}\lambda_j^{-1/2}\notag\\
 &\le\frac{2\beta\sqrt J}{D\sqrt c}=O(\varepsilon),
 \label{sgt:prep:first-loss}
\end{align}
because $D\ge J$ and \eqref{sgt:prep:cutoff} holds.
The renewal equation gives mean total filter time until acceptance
$\mathbb ET_{\rm attempt}/s$; time and success need not be independent
within a trial. The accepted state remains exact on the supported
spectrum. This supplies a mean, not a maximum, and assumes prompt
jump registration, controlled false aborts and finite resets. Input
support sorting and reset times are additional costs.

\subsection{The physical clock, locality, and finite reference support}
\label{sgt:prep:physical-scope}

Let $\Lambda=\max_{j\le J}\lambda_j$. With maximal implementable loss
rate $\Gamma_{\rm hw}$ use
$L_{\rm ph}=(\Gamma_{\rm hw}/\Lambda)L$, so that
\begin{equation}
 t_{\rm ph}=\frac{\Lambda}{\Gamma_{\rm hw}}u.
 \label{sgt:prep:physical-clock}
\end{equation}
The needed upper bound $\Lambda=O_\delta(\varepsilon^{-1})$
does not follow from the lower spectral bound
\eqref{sgt:prep:spectral-bound}. It follows instead from the actual
counting law, Eq.~\eqref{sgt:geo:counting}: its coefficient satisfies
$\mathcal N_\alpha(E)/E\to1/(2\pi)$ as $E\to\infty$.
Indeed the local matrix $S(q)$ tends to
$\operatorname{diag}(\alpha_i)$, its inverse-energy branches obey
$e_i(s,q)\sim2q/[w(s)\alpha_i]$, and their radial crossings satisfy
$q_i(s,E)/E\to w(s)\alpha_i/2$, uniformly in $s$.
Integrate and use the two angular signs and total mass one.
Choose a fixed $E$ with $\mathcal N_\alpha(E)>\kappa_\delta$,
and similarly for the base reference. The strict margin in the actual
limit $\varepsilon\#\{\lambda_j\le E/\varepsilon\}
\to\mathcal N_\alpha(E)$ pays the ceiling in $J$.
Thus \eqref{sgt:prep:duration-cutoff} gives a bounded maximum
physical filter time $O_{\delta,\delta_p,\tau}(\Gamma_{\rm hw}^{-1})$,
and \eqref{sgt:prep:first-loss} gives a bounded mean with monitoring.
This does not provide a numerical onset, a uniform simultaneous
$\delta(\varepsilon)\to0$ estimate, or a size-independent realizable
$\Gamma_{\rm hw}$.

If a mass-normalized graph stiffness actually has
\[
\begin{gathered}
 K=\sum_{ij}g_{ij}(e_i-e_j)(e_i-e_j)^*
       +\sum_i g_{i0}e_ie_i^*,\\
 L=M_b^{-1/2}KM_b^{-1/2},\\
 M_b=\operatorname{diag}(m_i)>0,
\end{gathered}
\]
the jumps
$\sqrt{g_{ij}}(a_i/\sqrt{m_i}-a_j/\sqrt{m_j})$ and
$\sqrt{g_{i0}/m_i}\,a_i$ implement that $L$.
Their neighboring modes require a common loss reservoir.
For the actual curve operator, eliminating the carrier gives a
Dirichlet-to-Neumann map, or on a mesh the Schur complement
\[
\begin{gathered}
 K_{\rm eff}=K_{BB}-K_{BI}K_{II}^{-1}K_{IB},\\
 L_{\rm eff}=M_B^{-1/2}K_{\rm eff}M_B^{-1/2}.
\end{gathered}
\]
This is generally dense even if the uneliminated carrier is local
\cite{DorflerBullo2013}. Static elimination does not establish that
ordinary local bulk damping produces $e^{-uL_{\rm eff}}$ on the
boundary. Positive bulk inertia adds dynamical modes and a
frequency-dependent elimination. A specimen construction must prove
its common mode map, generator approximation, heat and weighted-tail
errors, and resources across the growing band.

When $\lambda_j$ denotes squared mechanical frequency, a retained
single excitation has maximum frequency
$O_\delta(\varepsilon^{-1/2})$ and its energy, in fixed physical
units, must be supplied and reset. An optical path encoding at a
fixed photon carrier reproduces the encoded state; that carrier
energy is not the mechanical excitation energy. Native mechanical
oscillators with different frequency matrices generally have
different annihilation operators and vacua, related by a Bogoliubov
transformation. Second quantization of the mathematical $L^2$
alignment specifies a common number-conserving Fock model; it does
not prove a passive, vacuum-preserving laboratory identification.

A valid continuum soft reference retaining longitudinal coercivity is
\[
 H_{\rm soft}=L_{\rm base}^{\gamma}\otimes I_M+\Delta I\otimes Q.
\]
Its one-excitation Gibbs weight on $Q$ and its trace distance from the
hard reference are both
\begin{equation}
 p_{Q,\rm soft}=
 \frac{(M-1)e^{-\beta\Delta}}
 {1+(M-1)e^{-\beta\Delta}}.
 \label{sgt:prep:soft-reference}
\end{equation}
Hence error at most $\delta_{\rm ref}$ follows from
$\beta\Delta\ge
\log[(M-1)(1-\delta_{\rm ref})/\delta_{\rm ref}]$ when this is positive.
A constant penalty on an otherwise dispersionless infinite $Q$ space
would have infinite thermal degeneracy. Also a spring changes squared
frequencies; it does not automatically realize the displayed additive
frequency-Hamiltonian penalty. Full Bose vacuum-reference accuracy
has a separate population budget, as discussed with the Bose model.

\paragraph{Error and resource bookkeeping.}
If the continuum truncation distance is $\delta$, the additional
preparation distance is $\delta_p$, and the effect error is at most
$\delta_Q$ in operator norm, the measured $Q$ probability differs
from its continuum value by at most $\delta+\delta_p+\delta_Q$.
For two contractions $F,\widetilde F$ with
$\|\widetilde F-F\|\le d_F$,
\[
 \left\|\frac{\widetilde F\widetilde F^*-FF^*}{D}\right\|_1
 \le2d_F,
 \qquad D_{\rm tr}(\widetilde\rho,\rho)\le\frac{2d_F}{s},
\]
provided the perturbed successful output is nonzero. The first bound
follows by expanding the difference into two products and using
$\|X\|_1\le D\|X\|$; the second follows by comparing normalizations.
This is a whole-device norm tolerance, not a constant allowance on
each of $O(D^2)$ optical elements. Mode-dependent damping changes
the state. Even an extra scalar damping rate reweights the random
durations in the pooled subordination experiment and cannot generally
be removed as one overall transmission number. A sufficient small-loss
condition over a physical shot of duration $T$ is
$\max_j\kappa_jT\le\delta_{\rm loss}$; for
$\kappa_j=\omega_j/Q_{{\rm mech},j}$ this requires
$Q_{{\rm mech},j}\ge\omega_jT/\delta_{\rm loss}$.
Preparation and measured-copy cost therefore includes geometry and
coefficient computation, compilation and calibration, followed by
\[
 \frac{N}{s q_{\rm src}q_{\rm col}}
 \bigl(\text{injection}+
       \text{filter time}+\text{reset}+\text{readout}\bigr),
\]
with a renewal expectation for variable-duration trials and actual
conditional transmission where it is state dependent. Storage,
number resolution, dark counts, additional mode conversion and
control updates are separate resources. Only the ideal acceptance
floor and the displayed conditional filter-time estimates are proved
here.

% Standalone proof content for the SG-T manuscript.  Included after the
% geometric and controlled-preparation appendices; no private-note inputs.
\section{Full Bose ensembles: counting, thermodynamics and receiver costs}
\label{sgt:bose:appendix}

The normalized one-excitation heat state and the full bosonic equilibrium
state are different statistical experiments.  This appendix constructs the
second experiment in a specified common number-conserving representation,
proves its counting and thermodynamic laws, and identifies the resources
behind its discrimination bounds.  Its physical premise is a common
annihilation-mode encoding.  The geometric $L^2$ identification alone does
not establish that premise.

\subsection{The ensemble and the spectral estimates used below}
\label{sgt:bose:model}

Fix the curve, carrier and permitted distant masks, the smooth positive
density $w$ with $\int_\Gamma w\,\dd s=1$, and a finite profile with $M\ge2$
distinct heights $a_j$ and weights $\alpha_j>0$, $\sum_j\alpha_j=1$.
All these data, the exponent $\gamma>0$ and $\tau>0$ are fixed before
$h\downarrow0$.  Set
\begin{equation}
\begin{gathered}
 \mathfrak h=L^2(\Gamma,\dd s)\otimes\mathbb C^M,
 \\
 v=(\sqrt{\alpha_j})_{j=1}^M,\\
 P=vv^*,\\
 Q=I-P,\\
 \beta=\tau h^\gamma.
 \label{sgt:bose:parameters}
\end{gathered}
\end{equation}
We also use $P,Q$ for their tensor extensions to $\mathfrak h$.
Write $R_{1,h}=R_h$ for the actual positive Green operator and
$L_{1,h}=R_h^{-1}$ for its positive mass eigenvalue operator.  The
collapsed physical Green operator is $R_{0,h}=\widetilde R$ on
$L^2(\Gamma)$, independent of $h$, and
$L_{0,h}=L_0=\widetilde R^{-1}$.  The map $f\mapsto f\otimes v$
identifies this scalar physical space with $P\mathfrak h$.  An inverse power of the aligned collapsed operator
is used only through its heat, extended by zero on $Q\mathfrak h$.
In particular define
\begin{equation}
\begin{gathered}
 \Theta_{1,h}=A_h=\ee^{-\beta L_{1,h}^{\gamma}},\\
 \Theta_{0,h}=A_{0,h}=\ee^{-\beta L_0^{\gamma}}\otimes P,
 \\
 N_{i,h}=\Theta_{i,h}(I-\Theta_{i,h})^{-1}.
 \label{sgt:bose:heat-covariance}
\end{gathered}
\end{equation}
Thus $\Theta_{1,h},\Theta_{0,h}$ are local aliases for the heat
operators $A_h,A_{0,h}$, whereas the scaled Green operator in the
geometric appendix is $X_h=R_h/h$.  The index $i=0$ always refers to the
physical collapsed spectrum with a vacuum extension on $Q$.

The uniform spectral bound proved in
\eqref{sgt:prep:spectral-bound} gives constants $c,C>0$, independent of sufficiently
small $h$, such that the nonzero physical eigenvalues satisfy
\begin{equation}
\begin{gathered}
 r_j(R_{i,h})\le Cj^{-1},\\
 \lambda_j(L_{i,h})\ge cj,\\
 j\ge1,\\
 i=0,1.
 \label{sgt:bose:spectral-input}
\end{gathered}
\end{equation}
There are no additional zero-energy eigenvalues on the reference
complement.  Recall why the estimate is uniform: each diagonal self-curve
Green block is a weighted periodic logarithmic operator plus a uniformly
smooth remainder and maps $L^2$ to $H^1$ uniformly.  Its singular values
are $O(j^{-1})$.  Positivity of the full block operator gives
$0\le R_h\le M\operatorname{diag}(R_{11,h},\ldots,R_{MM,h})$ by applying
$\|\sum_{j=1}^M y_j\|^2\le M\sum_j\|y_j\|^2$ to
$y_j=R_h^{1/2}I_jx_j$.  Min--max and the fixed finite direct sum give
\eqref{sgt:bose:spectral-input}.  Thus this argument never requires
off-diagonal collision kernels to remain uniformly smooth.

For any $t>0$, decreasing-integral comparison yields
\begin{gather}
 \Tr\ee^{-tL_{i,h}^{\gamma}}
 \le\sum_{j\ge1}\ee^{-t(cj)^\gamma}
 \le\int_0^\infty\ee^{-t(cx)^\gamma}\,\dd x
 \le K_\gamma t^{-1/\gamma},
 \label{sgt:bose:heat-bound}\\
 h\Tr\!\left[(k\beta L_{i,h}^{\gamma})^p
                 \ee^{-k\beta L_{i,h}^{\gamma}}\right]
 \le K_{p,\gamma}(k\tau)^{-1/\gamma},
 \\
 k\ge1,\\
 p=0,1,2.
 \label{sgt:bose:weighted-bound}
\end{gather}
The second estimate follows for $p>0$ from
$x^p\ee^{-x}\le c_p\ee^{-x/2}$; all reference traces are on its
physical spectrum.  For each fixed $h>0$, $\Theta_{i,h}$ is positive
trace class with norm less than one.  Therefore $N_{i,h}$ is positive
trace class.  The same lower spectral bound implies
\begin{equation}
 \|N_{i,h}\|\le C h^{-\gamma},\qquad
 \nu_j(N_{i,h})\le C h^{-\gamma}j^{-\gamma},
 \label{sgt:bose:occupation-bound}
\end{equation}
where $\nu_j$ are the decreasing nonzero eigenvalues.  These estimates,
not the bounded heat-test theorem alone, control the Bose singularity
at small dimensionless energy.

For a contraction $T$ on $\mathfrak h$, let
$\Gamma_{\mathrm s}(T)=\bigoplus_{n\ge0}T^{\otimes_{\mathrm s}n}$,
with value one on the vacuum sector.  On
$\mathcal F_{\mathrm s}(\mathfrak h)$, define
\begin{equation}
 \mathcal Z_{i,h}=\det(I-\Theta_{i,h})^{-1},\qquad
 \varrho_{i,h}=\det(I-\Theta_{i,h})\Gamma_{\mathrm s}(\Theta_{i,h}).
 \label{sgt:bose:state}
\end{equation}
Diagonalizing $\Theta$ with eigenvalues $0\le\theta_j<1$ gives
$\Tr_{\mathcal F}\Gamma_{\mathrm s}(\Theta)=\prod_j(1-\theta_j)^{-1}$.
The product is finite because
$\sum_j-\log(1-\theta_j)\le\Tr\Theta/(1-\|\Theta\|)$.
This proves existence and normalization, including the infinite number
of physical modes.  The signal has excitation Hamiltonian
$\mathcal H_{1,h}=\mathrm d\Gamma(L_{1,h}^{\gamma})$.
The reference is the Gibbs state of $\mathrm d\Gamma(L_0^\gamma)$ on
$\mathcal F_{\mathrm s}(P\mathfrak h)$, tensored with exact vacuum on
$\mathcal F_{\mathrm s}(Q\mathfrak h)$.  It can be described as exclusion
of that sector at infinite energy; it is not a finite-energy Gibbs
Hamiltonian on the whole common space.  All energies below are excitation
energies relative to these stipulated vacua; no zero-point or Casimir
energy comparison is included.  The isolated Hamiltonian conserves total
number, while preparation at zero chemical potential allows a bath to
exchange quanta.  A fixed-total-number canonical ensemble is a different
model.

\subsection{A noncommuting counting determinant}
\label{sgt:bose:determinant-section}

Let $E$ be any orthogonal projection on $\mathfrak h$ and
$0\le\eta\le1$ the independent detection efficiency.  The no-detection
effect is $\Gamma_{\mathrm s}(I-\eta E)$, since its value on $m$ quanta
in $E\mathfrak h$ is $(1-\eta)^m$.  At $\eta=1$ this means the vacuum
projector on that factor, including value one at $m=0$.
For either state in \eqref{sgt:bose:state}, or for any positive
trace-class $\Theta$ of norm less than one, one has
\begin{equation}
 p_{\mathrm{no}}(\eta,E)
 =\det_{E\mathfrak h}(I_E+\eta ENE)^{-1},
 \qquad N=\Theta(I-\Theta)^{-1}.
 \label{sgt:bose:no-click}
\end{equation}
In particular $E$ need not commute with $\Theta$, $N$ or the energy.

To prove the formula, put $T=I-\eta E$.  Multiplicativity of second
quantization and a positive cyclic sandwich give
\begin{align*}
 \Tr_{\mathcal F}\bigl[\Gamma_{\mathrm s}(\Theta)
                                \Gamma_{\mathrm s}(T)\bigr]
 &=\Tr_{\mathcal F}\Gamma_{\mathrm s}
                  (\Theta^{1/2}T\Theta^{1/2})\\
 &=\det(I-\Theta+\eta\Theta^{1/2}E\Theta^{1/2})^{-1}.
\end{align*}
The cyclic step is legitimate: $\Gamma_{\mathrm s}(\Theta^{1/2})$ is
Hilbert--Schmidt because its squared Hilbert--Schmidt norm equals
$\Tr_{\mathcal F}\Gamma_{\mathrm s}(\Theta)<\infty$, and
$\Gamma_{\mathrm s}(T)$ is bounded.  The sandwich is positive, trace
class and bounded above by $\Theta$, so the partition product applies.
Factor the one-particle determinant by $I-\Theta$ to obtain
$\det(I+\eta N^{1/2}EN^{1/2})^{-1}$ after normalization.
For completeness, the required determinant algebra can first be applied
to spectral truncations $\Theta_m$; then
$\Theta_m\to\Theta$ and $N_m\to N$ in trace norm, and all displayed
positive sandwiches converge in trace norm.  Continuity of the
Fredholm determinant passes the finite-dimensional identities to the
limit.  Finally $N^{1/2}E$ and $EN^{1/2}$ are Hilbert--Schmidt and the
identity $\det(I+XY)=\det(I+YX)$ moves the determinant to
$E\mathfrak h$.  This proves \eqref{sgt:bose:no-click}.  The
finite-dimensional Fock trace identity is classical
\cite{Klich2002}; the preceding argument supplies the needed trace-class
extension.

The same calculation gives the probability generating function for the
detected $E$-count $K$:
\begin{equation}
 \mathbb E z^K
   =\det_{E\mathfrak h}
       [I_E+\eta(1-z)ENE]^{-1},\qquad 0\le z\le1.
 \label{sgt:bose:pgf}
\end{equation}
It follows either by conditioning the independent loss on the number
in $E$, or by replacing $\eta$ in the no-click identity with
$\eta(1-z)$.

For $q=|\xi|>0$, retain the geometric symbol
\begin{equation}
\begin{gathered}
 B(s,\xi)=\frac{w(s)}{2q}
   [\sqrt{\alpha_i\alpha_j}\ee^{-q|a_i-a_j|}]_{i,j=1}^M,
 \\
 C_{Q,\gamma}(t)=\frac1{2\pi}\int_\Gamma\int_{\mathbb R}
            \operatorname{tr}(Q\ee^{-tB(s,\xi)^{-\gamma}})
                         \,\dd\xi\,\dd s.
 \label{sgt:bose:Q-coefficient}
\end{gathered}
\end{equation}
The heat law \eqref{sgt:geo:heat-law}, with its fixed matrix insertion,
gives $h\Tr(Q\Theta_{1,h})\to C_{Q,\gamma}(\tau)>0$.
Positivity is strict: $B(s,\xi)>0$ for $\xi\ne0$, and $Q$ is a
nonzero projection.  Since $N\ge\Theta$, positive compression,
eigenvalue monotonicity and
$\log(1+\eta x)\ge x\log(1+\eta)$ for $0\le x\le1$ show that
\begin{equation}
\begin{gathered}
 -\log p_{\mathrm{no}}(\eta,Q)
 =\Tr_Q\log(I_Q+\eta QNQ)
\\
 \ge\log(1+\eta)\Tr(Q\Theta).
 \label{sgt:bose:exponent}
\end{gathered}
\end{equation}
Consequently, for every fixed $0<\eta\le1$ and any
$0<c<\log(1+\eta)C_{Q,\gamma}(\tau)$, all sufficiently small $h$ satisfy
\begin{equation}
\begin{gathered}
 p_{\mathrm{no},1}(\eta,Q)\le\ee^{-c/h},\\
 P_{\mathrm{err}}^{\mathrm{click}}
   =\tfrac12p_{\mathrm{no},1}(\eta,Q)\le\tfrac12\ee^{-c/h},
 \\
 \tfrac12\|\varrho_{1,h}-\varrho_{0,h}\|_1
     \ge1-p_{\mathrm{no},1}(\eta,Q)\longrightarrow1.
 \label{sgt:bose:ideal-error}
\end{gathered}
\end{equation}
The equal-prior rule declares the signal if any transverse quantum is
detected.  The reference has zero false-positive probability because
$QN_{0,h}=0$.  The trace-distance lower bound follows by using this
same effect in the variational characterization of trace distance.
No claim of finite-$h$ Helstrom optimality is made.  For $k$ independent
whole preparations and the ideal reference, an any-click-across-gates
rule has error $p_{\mathrm{no},1}^k/2$; it consumes $k$ whole thermal
samples.

Finite Fourier detection suffices.  If $D_s=-\mathrm i\partial_s$ on
the periodic curve of length $\ell_\Gamma$ and $0<d<R<\infty$ are fixed,
let $E_h=\one_{\{d<|hD_s|<R\}}\otimes Q$.  Then
\begin{gather}
 \operatorname{rank}E_h
   =\frac{(M-1)\ell_\Gamma(R-d)}{\pi h}+O(1),
 \label{sgt:bose:band-rank}\\
 h\Tr(E_h\Theta_{1,h})
\\
   \longrightarrow\frac1{2\pi}
       \int_\Gamma\int_{d<|\xi|<R}
          \operatorname{tr}(Q\ee^{-\tau B^{-\gamma}})
                       \,\dd\xi\,\dd s>0.
 \label{sgt:bose:band-signal}
\end{gather}
The second assertion is the sharp Fourier insertion law
\eqref{sgt:geo:band-law}.  Its treatment of the discontinuous band edges
uses smooth multiplier insertions, then bounds the two transition strips
by their $O(\text{strip width}/h)$ ranks and $0\le\Theta\le I$.
Thus it concerns the actual geometric heat and its complete tails.
Applying \eqref{sgt:bose:exponent} to $E_h$ proves an exponential
error bound with finitely many, $O(h^{-1})$, measured modes.  This rank
count does not supply a bandwidth, a mode converter or a simultaneous
readout apparatus.

\subsection{All-mode Bose thermodynamics and the price of one sample}
\label{sgt:bose:thermodynamics}

For $p=0,1,2$ define the split coefficients
\begin{equation}
\begin{gathered}
 C_{p,\gamma}(t)=\frac1{2\pi}
  \int_\Gamma\int_{\mathbb R}
   \operatorname{tr}\bigl[(tB^{-\gamma})^p
                              \ee^{-tB^{-\gamma}}\bigr]
                             \,\dd\xi\,\dd s,
 \\
 C_\gamma=C_{0,\gamma}.
 \label{sgt:bose:weighted-coefficients}
\end{gathered}
\end{equation}
Reference coefficients will carry a superscript $\mathrm{col}$, to
distinguish a collapsed index from the moment order $p=0$.
The actual convergence to these integrals is
\eqref{sgt:geo:weighted-heat}.  Put
\begin{equation}
\begin{gathered}
 \mathcal L_\gamma(\tau)=\sum_{k\ge1}\frac{C_\gamma(k\tau)}k,
 \\
 \mathcal E_\gamma(\tau)=\sum_{k\ge1}\frac{C_{1,\gamma}(k\tau)}k,
 \\
 \mathcal V_\gamma(\tau)=\sum_{k\ge1}\frac{C_{2,\gamma}(k\tau)}k.
 \label{sgt:bose:scalar-coefficients}
\end{gathered}
\end{equation}
All three series converge for every fixed $\gamma>0$ and $\tau>0$.
With $\langle\cdot\rangle_{1,h}$ denoting expectation in the signal
state, the limits are
\begin{gather}
 h\log\mathcal Z_{1,h}\longrightarrow\mathcal L_\gamma(\tau),
 \\
h\beta\langle\mathcal H_{1,h}\rangle_{1,h}
    \longrightarrow\mathcal E_\gamma(\tau),\nonumber\\
 h\beta^2\operatorname{Var}_{1,h}(\mathcal H_{1,h})
    \longrightarrow\mathcal V_\gamma(\tau),
 \\
hS(\varrho_{1,h})
    \longrightarrow\mathcal L_\gamma(\tau)+\mathcal E_\gamma(\tau).
 \label{sgt:bose:scalar-laws}
\end{gather}
For each fixed $r>0$, $r\ne1$, the R\'enyi entropy satisfies
\begin{equation}
 hS_r(\varrho_{1,h})\longrightarrow
 \frac{\mathcal L_\gamma(r\tau)-r\mathcal L_\gamma(\tau)}{1-r}.
 \label{sgt:bose:renyi-law}
\end{equation}
The same statements hold for the collapsed physical ensemble.

Here are the convergent expansions that prove these claims.  Each
physical normal mode of energy $\omega$ has a geometric occupation of
mean $(\ee^{\beta\omega}-1)^{-1}$ and variance
$\ee^{\beta\omega}/(\ee^{\beta\omega}-1)^2$; different normal modes
are independent by the product state.  Expanding the geometric series
and using nonnegative sums gives
\begin{align}
 \log\mathcal Z_{1,h}
    &=\sum_{k\ge1}\frac1k\Tr\Theta_{1,h}^k,
 \nonumber\\
 \beta\langle\mathcal H_{1,h}\rangle_{1,h}
    &=\sum_{k\ge1}\frac1k
      \Tr\bigl[(k\beta L_{1,h}^{\gamma})
                   \ee^{-k\beta L_{1,h}^{\gamma}}\bigr],
 \nonumber\\
 \beta^2\operatorname{Var}_{1,h}(\mathcal H_{1,h})
    &=\sum_{k\ge1}\frac1k
      \Tr\bigl[(k\beta L_{1,h}^{\gamma})^2
                   \ee^{-k\beta L_{1,h}^{\gamma}}\bigr].
 \label{sgt:bose:series-proof}
\end{align}
After multiplication by $h$, \eqref{sgt:bose:weighted-bound} dominates
each term by a constant times $k^{-1-1/\gamma}$, a summable sequence for
every $\gamma>0$.  The weighted heat theorem therefore applies
term by term, proving the first three limits.  Finiteness at each $h$
and the exact Gibbs identity
$S=\log\mathcal Z+\beta\langle\mathcal H\rangle$ prove the entropy
limit.  Similarly
$\Tr\varrho_{1,h}^r=\mathcal Z_{1,h}(r\beta)/\mathcal Z_{1,h}(\beta)^r$
proves \eqref{sgt:bose:renyi-law}.  This argument does not differentiate
an uncontrolled asymptotic remainder or invoke entropy continuity in
unbounded dimension.

Occupation laws require the additional restriction $0<\gamma<1$.
In this range define
\begin{equation}
\begin{gathered}
 \mathsf n_1=\mathcal N_\gamma(\tau)=\sum_{k\ge1}C_\gamma(k\tau),
 \\
 \mathsf n_Q=\mathcal N_{Q,\gamma}(\tau)
      =\sum_{k\ge1}C_{Q,\gamma}(k\tau)>0.
 \label{sgt:bose:number-coefficients}
\end{gathered}
\end{equation}
Because $N_{1,h}=\sum_{k\ge1}\Theta_{1,h}^k$ in trace norm, the heat
bound now supplies the summable majorant $k^{-1/\gamma}$.  Positivity
also bounds a $Q$-inserted trace by the total trace.  Hence
\begin{equation}
\begin{gathered}
 h\Tr N_{1,h}\longrightarrow \mathsf n_1,
 \\
 h\Tr(QN_{1,h})\longrightarrow \mathsf n_Q,
 \\
 h\Tr N_{0,h}\longrightarrow \mathsf n_0
      :=\mathcal N^{\mathrm{col}}_\gamma(\tau).
 \label{sgt:bose:number-law}
\end{gathered}
\end{equation}
For $\gamma\ge1$ this dominating series is not summable, so no
$h^{-1}$ occupation law is asserted here.  This restriction does not
affect \eqref{sgt:bose:scalar-laws}--\eqref{sgt:bose:renyi-law}.
The cost of a whole preparation, for $0<\gamma<1$, is consequently
\begin{equation}
 \langle\mathrm d\Gamma(I)\rangle_{1,h}
     =\frac{\mathsf n_1+o(1)}h,
 \qquad
 \langle\mathcal H_{1,h}\rangle_{1,h}
     =\frac{\mathcal E_\gamma(\tau)+o(1)}{\tau h^{1+\gamma}}.
 \label{sgt:bose:resources}
\end{equation}
In particular the frequency exponent $\gamma=1/2$ requires
$O(h^{-1})$ mean quanta and $O(h^{-3/2})$ mean excitation energy, with
strictly positive leading coefficients.  Equilibration, source and
reset times, dissipation, and conversion work are additional resources.

Conditioning \eqref{sgt:bose:state} on exactly one excitation gives
the normalized one-particle heat state.  Its exact probability is
\begin{equation}
\begin{gathered}
 p_{1,h}=\frac{\Tr\Theta_{1,h}}{\mathcal Z_{1,h}},
 \\
 p_{1,h}\le(\Tr\Theta_{1,h})\ee^{-\Tr\Theta_{1,h}},
 \\
 h\log p_{1,h}\longrightarrow-\mathcal L_\gamma(\tau)<0.
 \label{sgt:bose:postselection}
\end{gathered}
\end{equation}
The inequality follows from $-\log(1-x)\ge x$.  For the last limit,
$h\Tr\Theta_{1,h}\to C_\gamma(\tau)>0$ implies
$h\log\Tr\Theta_{1,h}\to0$, while the partition limit is already
proved.  Thus the full-ensemble receiver avoids an exponentially rare
conditioning event by using a different sample, with the extensive
cost in \eqref{sgt:bose:resources}.

The reference coefficients can be evaluated explicitly.  Put
$\sigma=1/\gamma$ and $d_\gamma=\Gamma(1+\sigma)/(2\pi)$.
Integration of the rank-one reference symbol gives
$C_\gamma^{\mathrm{col}}(t)=d_\gamma t^{-\sigma}$,
$C_{1,\gamma}^{\mathrm{col}}=\sigma C_\gamma^{\mathrm{col}}$ and
$C_{2,\gamma}^{\mathrm{col}}=\sigma(\sigma+1)C_\gamma^{\mathrm{col}}$.
The corresponding Bose constants are
\begin{gather}
 \mathcal L_\gamma^{\mathrm{col}}
  =d_\gamma\zeta(1+\sigma)\tau^{-\sigma},
 \\
\mathcal E_\gamma^{\mathrm{col}}
  =\sigma d_\gamma\zeta(1+\sigma)\tau^{-\sigma},\nonumber\\
 \mathcal V_\gamma^{\mathrm{col}}
  =\sigma(\sigma+1)d_\gamma\zeta(1+\sigma)\tau^{-\sigma},
 \\
\mathcal N_\gamma^{\mathrm{col}}
  =d_\gamma\zeta(\sigma)\tau^{-\sigma}\\
(\sigma>1).
 \label{sgt:bose:reference-coefficients}
\end{gather}
The zeta function here is simply $\zeta(s)=\sum_{k\ge1}k^{-s}$ on
its convergent real range.

The high-resolution endpoints of the geometric theorem,
\eqref{sgt:geo:hot-endpoint} and \eqref{sgt:geo:moment-endpoint},
extend to these convergent sums:
\begin{equation}
 \frac{\mathcal L_\gamma}{\mathcal L_\gamma^{\mathrm{col}}},\quad
 \frac{\mathcal E_\gamma}{\mathcal E_\gamma^{\mathrm{col}}},\quad
 \frac{\mathcal V_\gamma}{\mathcal V_\gamma^{\mathrm{col}}}
 \longrightarrow1\qquad(\tau\downarrow0),
 \label{sgt:bose:hot-scalars}
\end{equation}
and, for $0<\gamma<1$,
\begin{equation}
 \frac{\mathcal N_\gamma}{\mathcal N_\gamma^{\mathrm{col}}}
   \longrightarrow1,
 \qquad
 \frac{\mathcal N_{Q,\gamma}}{\mathcal N_\gamma}
   \longrightarrow1-\sum_{j=1}^M\alpha_j^2.
 \label{sgt:bose:hot-numbers}
\end{equation}
To check the interchange in this second, successive limit, the eigenvalues
of the transverse matrix in $B$ are at most its trace one.  With
$\xi=w(s)t^{-1/\gamma}u$, this gives
$C_{p,\gamma}(t)\le c_p t^{-1/\gamma}$ for all $t>0$, using
$x^p\ee^{-x}\le c_p\ee^{-x/2}$ if necessary.  Thus
$\tau^{1/\gamma}C_{p,\gamma}(k\tau)$ is bounded by a constant times
$k^{-1/\gamma}$, independently of $\tau$.  The already proved hot
endpoint for each fixed $k$ and dominated convergence prove
\eqref{sgt:bose:hot-scalars}; for numbers the same argument is summable
exactly when $\gamma<1$.  The $Q$ endpoint has the extra factor
$1-\sum_j\alpha_j^2$, proving \eqref{sgt:bose:hot-numbers}.
It follows in particular that the von Neumann entropy coefficient and
each fixed positive-order R\'enyi entropy coefficient have ratio one
to their collapsed counterparts.  For $r\ne1$ the latter counterpart is
\begin{equation}
 d_\gamma\zeta(1+\sigma)\tau^{-\sigma}
                  \frac{r^{-\sigma}-r}{1-r}>0;
 \label{sgt:bose:reference-renyi}
\end{equation}
the separate $r=1$ coefficient is
$(1+\sigma)d_\gamma\zeta(1+\sigma)\tau^{-\sigma}$.

These are ratios of extensive full-Fock coefficients.  They are not the
additive one-excitation entropy agreement.  Indeed the strict geometric
partition deficit \eqref{sgt:geo:partition-deficit} implies, term by term,
\begin{equation}
\begin{gathered}
 \Delta_\gamma(\tau)
   :=\mathcal L_\gamma^{\mathrm{col}}(\tau)
                  -\mathcal L_\gamma(\tau)>0,
 \\
 \frac{\mathcal Z_{1,h}}{\mathcal Z_{0,h}}
   =\exp\!\left[-\frac{\Delta_\gamma(\tau)+o(1)}h\right]
   \longrightarrow0
 \label{sgt:bose:partition-ratio}
\end{gathered}
\end{equation}
for every fixed $\tau>0$ and every nontrivial fixed profile.
No arbitrary simultaneous $h,\tau$ limit, no moving profile, and no
uniform limit in increasing $M$ is included in these conclusions.

\subsection{The efficiency threshold and its Poisson transition}
\label{sgt:bose:efficiency}

Assume $0<\gamma<1$ in this subsection.  For any efficiency sequence
$0\le\eta_h\le1$, \eqref{sgt:bose:exponent} and
$\log(1+x)\le x$ give
\begin{equation}
\begin{gathered}
 \log(1+\eta_h)\Tr(Q\Theta_{1,h})
\\
 \le-\log p_{\mathrm{no},1}(\eta_h,Q)
 \le\eta_h\Tr(QN_{1,h}).
 \label{sgt:bose:loss-sandwich}
\end{gathered}
\end{equation}
Since $\log(1+\eta_h)\ge\eta_h/2$, the heat and number laws imply
\begin{equation}
\begin{gathered}
 p_{\mathrm{no},1}(\eta_h,Q)\to0
       \\
\Longleftrightarrow\\
\eta_h/h\to\infty,
 \\
 \eta_h/h\to0\\
\Longrightarrow\\
 p_{\mathrm{no},1}\to1.
 \label{sgt:bose:loss-threshold}
\end{gathered}
\end{equation}
For necessity in the equivalence, failure of $\eta_h/h\to\infty$
produces a subsequence with bounded $\eta_h/h$, on which the upper bound
in \eqref{sgt:bose:loss-sandwich} keeps the no-click probability bounded
away from zero.  This is an exact consistency criterion for this specified
vacuum-reference receiver, not an optimality theorem for all measurements.

If $\eta_h/h\to\kappa\in[0,\infty)$, let
$X_h=\eta_h QN_{1,h}Q$.  Then
$\Tr X_h\to\kappa \mathsf n_Q$ and
$\|X_h\|=O(h^{1-\gamma})\to0$.  Spectral calculus and the scalar
bound $0\le x-\log(1+x)\le x^2/2$ give
\begin{equation}
 0\le\Tr X_h-\Tr\log(I+X_h)
   \le\tfrac12\|X_h\|\Tr X_h\longrightarrow0.
 \label{sgt:bose:poisson-remainder}
\end{equation}
Apply this bound to $(1-z)X_h$ in \eqref{sgt:bose:pgf}.  The generating
function tends to $\exp[-(1-z)\kappa \mathsf n_Q]$ on $0\le z\le1$.
Equivalently, evaluating at $z=\ee^{-t}$ gives convergence of the
nonnegative-variable Laplace transforms to a transform continuous at
$t=0$.  Hence
\begin{equation}
 K_h\ \Longrightarrow\ \operatorname{Poisson}(\kappa \mathsf n_Q),
 \qquad
 p_{\mathrm{no},1}(\eta_h,Q)\longrightarrow\ee^{-\kappa \mathsf n_Q}.
 \label{sgt:bose:poisson}
\end{equation}
The Poisson approximation is valid in this vanishing-efficiency regime.
It is not used for the highly occupied modes of a fixed-efficiency
thermal sample.

\subsection{Why an any-click decision is fragile}
\label{sgt:bose:fragility}

Suppose a whole detection gate has an independent dark-click probability
$d_h$, common to the hypotheses.  Write $p_1,p_0$ for their no-click
probabilities before this background.  The equal-prior any-click error
is exactly
\begin{equation}
 P_{\mathrm{err}}=\tfrac12
       \{1-(1-d_h)p_0+(1-d_h)p_1\}.
 \label{sgt:bose:dark-any-click}
\end{equation}
For the ideal reference $p_0=1$, the limiting floor is $d_h/2$ when
$p_1\to0$.  If a fixed $d<1$ applies independently to each of $k$
independent preparations, the rule declaring signal only when every
gate clicks has instead
\begin{equation}
 P_{\mathrm{err}}^{(k)}
  =\tfrac12\{d^k+1-[1-(1-d)p_1]^k\}
  \le\tfrac12\{d^k+k(1-d)p_1\}.
 \label{sgt:bose:dark-repeated}
\end{equation}
Thus a finite $k$ suffices for each fixed target error when $p_1$ is
small enough; it pays $k$ preparation costs.  When $D_h$ independent
channels each have dark-click probability $\delta_h$,
$d_h=1-(1-\delta_h)^{D_h}$.  For $D_h=O(h^{-1})$,
$\delta_h=o(h)$ suffices for a vanishing gate-level dark floor.
If instead $D_h\asymp h^{-1}$ and $\delta_h$ is a fixed positive
constant below one, $1-d_h$ is exponentially small in $h^{-1}$ and
making $d_h^k$ small can itself require exponentially many gates.
An aggregated detector has its own gate-level dark probability; the
independent-channel formula applies only when that is the actual model.

For a passive unitary sorter $V$, the detected projection is
$E=V^*QV$.  If $\theta_h=\|QVP\|$, the reference covariance satisfies
\begin{equation}
\begin{gathered}
 \Tr(EN_{0,h})=\Tr(PV^*QVPN_{0,h})
       \le\theta_h^2\Tr N_{0,h},\\
 \mathbb P_0(K\ge1)\le\eta\theta_h^2\Tr N_{0,h}.
 \label{sgt:bose:sorter-leakage}
\end{gathered}
\end{equation}
The last step is Markov's inequality.  For $0<\gamma<1$, keeping this
sufficient upper bound below a chosen small fixed $\delta$ requires
an amplitude budget of order
$\theta_h^2=O(\delta h/\eta)$, with the constant chosen from $\mathsf n_0$.
This is a worst-case sufficient bound, not a necessary bound for each
sorter.  The danger from fixed mixing is real: a fixed rotation from
$v$ into one fixed unit vector of $Q\mathbb C^M$ at every tangential
coordinate produces the reference output covariance
$\theta^2N_{0,h}$ on the rotated output channel.  For $\theta>0$ fixed,
\eqref{sgt:bose:exponent}, applied to the physical reference heat with
effective efficiency $\eta\theta^2$, makes its no-click probability
tend to zero exponentially.  Thus an extensive bright reference can
click almost surely under fixed coherent leakage.

A finite-energy reference has a separate quantitative budget.  Assume
it is defined in the same annihilation representation, has a positive
trace-class thermal covariance $N_{\mathrm r,h}$, commutes with $Q$,
and its one-particle energy restricted to $Q$ is at least
$\Delta_h I_Q$.  For a detected projection $E_h\le Q$ of finite rank
$D_h$, functional calculus gives
\begin{equation}
 1-p_{\mathrm r}\le\eta_h\Tr(E_hN_{\mathrm r,h})
   \le\frac{\eta_hD_h}{\ee^{\beta\Delta_h}-1}.
 \label{sgt:bose:reference-gap}
\end{equation}
Consequently the signal bound and \eqref{sgt:bose:dark-any-click} imply
\begin{equation}
\begin{gathered}
 P_{\mathrm{err}}
 \le\frac12\bigl\{d_h+
       \frac{\eta_hD_h}{\ee^{\beta\Delta_h}-1}
\\
       +\exp[-\log(1+\eta_h)\Tr(E_h\Theta_{1,h})]\bigr\}.
 \label{sgt:bose:reference-error}
\end{gathered}
\end{equation}
A spectral gap alone does not make an infinite-dimensional thermal
trace finite; existence of this reference and the finite detector rank
are explicit assumptions.  At fixed efficiency and $D_h=O(h^{-1})$,
the condition $\beta\Delta_h-\log D_h\to+\infty$ is sufficient for
a vanishing reference term.  Matching an $\ee^{-c/h}$ target by this
bound requires a corresponding additional $c/h$ in
$\beta\Delta_h$ and an equally controlled dark term.  These are
requirements of the vacuum-null decision rule.  The count threshold
below uses a calibrated, possibly nonzero, reference mean.

\subsection{Passive output covariances and a Bose concentration bound}
\label{sgt:bose:concentration-section}

We first justify the statistics after a common imperfect passive
receiver.  A gauge-invariant thermal state with positive trace-class
covariance $N$ has the exact coherent-state representation
\begin{equation}
 \varrho_N=\mathbb E\,|Z\rangle\langle Z|,
 \qquad \mathbb E Z=0,\qquad
 \mathbb E[ZZ^*]=N,
 \label{sgt:bose:gaussian-mixture}
\end{equation}
where $Z$ is a circular complex Gaussian vector in $\mathfrak h$ and
$|Z\rangle$ is the normalized coherent vector.  This assertion also
has an elementary infinite-dimensional justification.  In an eigenbasis
of $N$, put $Z_j=\sqrt{n_j}\xi_j$ with independent standard circular
complex Gaussians $\mathbb E|\xi_j|^2=1$.  Since $\sum_jn_j<\infty$,
Tonelli gives $\sum_j|Z_j|^2<\infty$ almost surely, so $Z$ lies in
$\mathfrak h$.  The integral of the positive rank-one projectors is a
trace-one operator.  In one mode, angular integration annihilates
off-diagonal number matrix elements and radial integration gives
\begin{equation}
 \int_{\mathbb C}\frac{\ee^{-|z|^2/n}}{\pi n}
       |z\rangle\langle z|\,\dd^2z
   =\sum_{m\ge0}\frac1{1+n}
                \left(\frac{n}{1+n}\right)^m|m\rangle\langle m|.
 \label{sgt:bose:one-mode-mixture}
\end{equation}
At $n=0$ use the vacuum delta measure.  Products give the equality on
every finite set of modes; increasing finite sets and the finite mean
total number identify the full trace-one states.  Taking
$n_j=\theta_j/(1-\theta_j)$ proves \eqref{sgt:bose:gaussian-mixture}
for \eqref{sgt:bose:state}.

A contraction $T:\mathfrak h\to\mathfrak g$, implemented by passive
mixing with vacuum ancillas followed by discarding unused outputs,
sends a coherent vector to a coherent vector of amplitude $TZ$ at
the measured output.  Equation \eqref{sgt:bose:gaussian-mixture}
therefore gives the output covariance $TNT^*$.  With a separate
efficiency $\eta$, it is
\begin{equation}
 C=\eta TNT^*.
 \label{sgt:bose:passive-covariance}
\end{equation}
For a detection effect $0\le E\le I$, take $T=\sqrt E$.  Diagonalizing
$C$ is a mathematical step only: total number is invariant under
passive output rotations, so the receiver need not know which
hypothesis-specific eigenbasis diagonalizes $C$.  Its total count is
a sum of independent geometric variables whose means are the
eigenvalues $c_j$ of $C$.

Set $\mu=\Tr C$ and $B_C=\|C\|$.  For
$|t|\le[4(1+B_C)]^{-1}$, the logarithmic moment generating function is
\begin{align}
 F(t)&=-\sum_j\log[1-c_j(\ee^t-1)],\nonumber\\
 F''(t)&=\sum_j
       \frac{c_j(1+c_j)\ee^t}{[1-c_j(\ee^t-1)]^2}.
 \label{sgt:bose:mgf}
\end{align}
For $0\le t\le[4(1+B_C)]^{-1}$ one has
$\ee^t\le4/3$ and $\ee^t-1\le(4/3)t$; every denominator before
squaring is at least $2/3$.  Thus
$F''(t)\le3\mu(1+B_C)$.  For negative $t$ each denominator is at
least one and the same bound follows.  These estimates also justify
termwise differentiation on compact interior intervals by
$\sum_jc_j<\infty$; continuity gives the endpoints.  Since
$F(0)=0$ and $F'(0)=\mu$, integration twice proves
\begin{equation}
 \log\mathbb E\ee^{t(K-\mu)}
       \le2\mu(1+B_C)t^2,
 \qquad |t|\le\frac1{4(1+B_C)}.
 \label{sgt:bose:centered-mgf}
\end{equation}
For either tail choose
$|t|=u/[4(1+B_C)(\mu+u)]$ in Chernoff's inequality.  Substitution
gives, for every $u>0$,
\begin{equation}
\begin{gathered}
 \mathbb P(K-\mu\ge u),\\
\mathbb P(K-\mu\le-u)
   \le\exp\!\left[-\frac{u^2}
                         {8(1+B_C)(\mu+u)}\right].
 \label{sgt:bose:concentration}
\end{gathered}
\end{equation}
If $\mu=0$, $K=0$ almost surely and the same tail bounds hold.

If $D$ is an independent Poisson dark count of known mean $\nu$, then
$\log\mathbb E\ee^{t(D-\nu)}=\nu(\ee^t-1-t)\le\nu t^2$ for
$|t|\le1/4$.  Combining with \eqref{sgt:bose:centered-mgf} proves
\eqref{sgt:bose:concentration} for $K+D$, centered at $\mu+\nu$,
with $\mu$ in the denominator replaced by $\mu+\nu$ and $B_C$
unchanged.  This is an independent calibrated Poisson allowance;
it is not a claim for arbitrary correlated noise or unknown drift.

\subsection{A count threshold tolerates fixed small effect error}
\label{sgt:bose:unfiltered}

Fix $0<\gamma<1$, $\tau>0$ and $0<\eta\le1$.  The constants
$\mathsf n_1,\mathsf n_0,\mathsf n_Q$ are those in \eqref{sgt:bose:number-law}.
Suppose the same effect $E_h$ is used under both hypotheses and
\begin{equation}
\begin{gathered}
 0\le E_h\le I,\\
 \|E_h-Q\|\le e,
 \\
 g:=\eta[\mathsf n_Q-e(\mathsf n_1+\mathsf n_0)]>0,
 \label{sgt:bose:effect-hypothesis}
\end{gathered}
\end{equation}
where $e$ is fixed.  The output covariances are
$C_{i,h}=\eta\sqrt{E_h}N_{i,h}\sqrt{E_h}$, with means
$\mu_{i,h}=\Tr C_{i,h}$.  Since $QN_{0,h}=0$, trace duality gives
\begin{gather}
 \mu_{1,h}-\mu_{0,h}
 \ge\eta\{\Tr(QN_{1,h})-e[\Tr N_{1,h}+\Tr N_{0,h}]\}
\\
   =\frac{g+o(1)}h,\nonumber\\
 \mu_{i,h}=O(h^{-1}),\\
 \|C_{i,h}\|=O(h^{-\gamma}).
 \label{sgt:bose:mean-gap}
\end{gather}
Take the threshold nearest to
$(\mu_{1,h}+\mu_{0,h})/2$, declaring signal for counts above it.
The two error margins are at least
$(\mu_{1,h}-\mu_{0,h})/2-1\ge c_1/h$ for all sufficiently small $h$.
Applying \eqref{sgt:bose:concentration} to both hypotheses yields a
constant $c_2>0$ such that
\begin{equation}
 P_{\mathrm{err}}^{\mathrm{threshold}}
      \le\exp[-c_2h^{\gamma-1}].
 \label{sgt:bose:unfiltered-rate}
\end{equation}
Thus the frequency exponent $\gamma=1/2$ gives
$\exp(-c_2/\sqrt h)$.  The occupation norm in
\eqref{sgt:bose:mean-gap} retains Bose bunching; replacing those
fluctuations by Poisson ones at fixed efficiency would not prove the
stated result.

An independent common Poisson dark count with known mean
$\nu_h\le C_\nu/h$ is included by adding $\nu_h$ to the threshold and
using the preceding concentration result.  It changes constants, not
the rate.  More generally a threshold is valid whenever known background
means, or rigorous calibration intervals for them, leave a positive
$h^{-1}$ separation and obey the same moment bound.  This condition
must be checked; arbitrary unknown offsets are not covered.  Total-count
readout must resolve counts over a range of order $h^{-1}$; a saturated
one-bit detector does not implement this decision.

A reference need not be nearly vacuum for such a threshold.  In a
finite detected sector of rank at most $d/h$, a reference occupation
bounded by $n_\Delta=(\ee^{\beta\Delta_h}-1)^{-1}$ contributes mean
at most $\eta d n_\Delta/h$.  It can be included in a calibrated null
provided the mean separation and covariance bound remain valid.  A
sufficiently large fixed $\beta\Delta_h$ may then suffice.  On an
infinite sector a constant gap without dispersion is not a trace-class
thermal state.  If instead $H_Q=H_{\mathrm{long}}+\Delta_h I$ with
positive confining $H_{\mathrm{long}}$, then the spectral inequality
\begin{equation}
 \Tr(\ee^{\beta H_Q}-1)^{-1}
  \le\frac{\ee^{-\beta\Delta_h}}{1-\ee^{-\beta\Delta_h}}
                         \Tr\ee^{-\beta H_{\mathrm{long}}}
 \label{sgt:bose:dispersive-reference}
\end{equation}
follows by expanding its geometric series, or by bounding the
denominator mode by mode.  A controlled $O(h^{-1})$ heat trace thus
supplies a finite reference number budget.  This is a model of an
additive frequency-energy gap, not an assertion that a stiffness
penalty adds linearly after taking a square root.

\subsection{One common low-frequency rejection space}
\label{sgt:bose:filter}

The previous rate is limited by the covariance of low-energy Bose
modes.  We now construct one common filter, designed from the two
specified candidate operators, that rejects those modes without
discarding the extensive transverse signal.  Continue to assume
$0<\gamma<1$ and fixed $\tau,\eta>0$.

For a fixed $b>0$, let
$F_{1,h}=\one_{[0,b/h)}(L_{1,h})$.  Let $F_{0,h}$ be the analogous
projection for $L_0$ on $P\mathfrak h$, extended by zero on $Q$.
By \eqref{sgt:bose:spectral-input},
$\operatorname{rank}F_{i,h}\le Cb/h$.  Define
\begin{equation}
\begin{gathered}
 W_h=\operatorname{ran}F_{1,h}+\operatorname{ran}F_{0,h}
                          +P\operatorname{ran}F_{1,h},
 \\
 r_h=\dim W_h\le C_Wb/h,
 \\
\Pi_h=I-\operatorname{proj}W_h.
 \label{sgt:bose:rejected-space}
\end{gathered}
\end{equation}
The space $W_h$ is invariant under $P$: $P$ maps its first summand
into its third, fixes its second, and fixes the image of $P$.
As $P=P^*$, its orthogonal complement is invariant too.  Hence
$[\Pi_h,P]=[\Pi_h,Q]=0$ and $\Pi_h\le I-F_{i,h}$ for both candidates.
Every vector in its range belongs to the high-energy spectral subspace
of each candidate, even though these two spectral projections need
not commute.  Functional calculus for each one separately gives
\begin{equation}
 \Pi_hN_{i,h}\Pi_h\le n_b\Pi_h,
 \qquad n_b=(\ee^{\tau b^\gamma}-1)^{-1},\quad i=0,1.
 \label{sgt:bose:bounded-filter-covariance}
\end{equation}
For the reference, the additional $Q$ component simply has zero
covariance.

For any rank-$r$ orthogonal projection $J$, Ky Fan's variational
inequality and \eqref{sgt:bose:occupation-bound} give
\begin{equation}
 \Tr(JN_{1,h})\le\sum_{j=1}^r\nu_j(N_{1,h})
   \le\frac{C}{1-\gamma}h^{-\gamma}r^{1-\gamma}.
 \label{sgt:bose:ky-fan}
\end{equation}
To see the inequality without any commutation hypothesis, diagonalize
$N$ and write $\Tr(JN)=\sum_j\nu_j\langle e_j,Je_j\rangle$; its
coefficients lie in $[0,1]$ and sum to $r$, so their largest possible
weighted sum fills the first $r$ entries.  The remaining bound uses
$\sum_{j=1}^rj^{-\gamma}\le r^{1-\gamma}/(1-\gamma)$.
Because $Q(I-\Pi_h)$ is itself an orthogonal projection of rank at
most $r_h$, this implies
\begin{equation}
\begin{gathered}
 0\le\Tr(QN_{1,h})-\Tr(Q\Pi_hN_{1,h}\Pi_h)
\\
   =\Tr[Q(I-\Pi_h)N_{1,h}]
\\
   \le C_*b^{1-\gamma}/h.
 \label{sgt:bose:filter-loss}
\end{gathered}
\end{equation}
The constant $C_*$ is independent of $b$ and sufficiently small $h$,
with the fixed geometric data, $\gamma$ and $\tau$ understood.
Choosing $b$ small therefore preserves a positive fraction of the
$\mathsf n_Q/h$ signal.  The choice is not uniform as $\gamma\uparrow1$.

A finite Fourier realization requires an intersection of spaces,
not a product of noncommuting projections.  Fix $R>0$ and let
$T_{R,h}=\one_{\{|hD_s|\le R\}}\otimes I_M$.  Let $J_{R,h}$ be the
orthogonal projection onto the finite-dimensional space $T_{R,h}W_h$,
and set
\begin{equation}
 \Xi_h=T_{R,h}-J_{R,h}.
 \label{sgt:bose:intersection-definition}
\end{equation}
The range of $J_{R,h}$ is contained in that of $T_{R,h}$, so $\Xi_h$
is an orthogonal projection.  For $x\in\operatorname{ran}T_{R,h}$,
$\langle x,T_{R,h}w\rangle=\langle x,w\rangle$ for every $w\in W_h$.
Consequently
\begin{equation}
\begin{gathered}
 \operatorname{ran}\Xi_h
   =\operatorname{ran}T_{R,h}\cap W_h^\perp,
 \\
\Xi_h\le\Pi_h,
 \\
[\Xi_h,P]=[\Xi_h,Q]=0.
 \label{sgt:bose:intersection}
\end{gathered}
\end{equation}
The commutation follows because $T_{R,h}$ commutes with $P$ and
$W_h$ is $P$-invariant, hence so is $T_{R,h}W_h$ and its orthogonal
projection.  In particular
\begin{equation}
 \operatorname{rank}\Xi_h
   \le\frac{M\ell_\Gamma R}{\pi h}+O(1),
 \qquad \Xi_hN_{i,h}\Xi_h\le n_b\Xi_h.
 \label{sgt:bose:finite-filter-bound}
\end{equation}
No commutation of $T_{R,h}$ with $\Pi_h$ or of $\Xi_h$ with either
Hamiltonian was used.

Since $I-\Xi_h=(I-T_{R,h})+J_{R,h}$ and all these projections commute
with $Q$, the exact loss of $Q$-number splits into two nonnegative
terms.  The rank of $J_{R,h}$ is at most $r_h$, so
\begin{equation}
\begin{gathered}
 \Tr(Q\Xi_hN_{1,h}\Xi_h)
\\
  \ge\Tr(QN_{1,h})-\Tr[Q(I-T_{R,h})N_{1,h}]
                         -C_*b^{1-\gamma}/h.
 \label{sgt:bose:finite-retention}
\end{gathered}
\end{equation}
The remaining high-frequency occupation tail is controlled on the
actual geometry.  Expand $N_{1,h}=\sum_{k\ge1}\Theta_{1,h}^k$, apply
the Fourier insertion law \eqref{sgt:geo:band-law} at each fixed
$k\tau$, and then use \eqref{sgt:bose:heat-bound} with the summable
$k^{-1/\gamma}$ majorant.  This proves
\begin{gather}
 h\Tr[Q(I-T_{R,h})N_{1,h}]\longrightarrow\mathcal T_Q(R),\nonumber\\
 \mathcal T_Q(R)
 =\frac1{2\pi}\int_\Gamma\int_{|\xi|>R}
    \sum_{k\ge1}\operatorname{tr}
        (Q\ee^{-k\tau B(s,\xi)^{-\gamma}})
                                      \,\dd\xi\,\dd s
   \longrightarrow0\\
(R\to\infty).
 \label{sgt:bose:fourier-number-tail}
\end{gather}
The final limit follows from positivity and finiteness of the full
integral $\mathsf n_Q$.  For example choose $b>0$ with
$C_*b^{1-\gamma}<\mathsf n_Q/8$ and then $R$ with
$\mathcal T_Q(R)<\mathsf n_Q/8$.  These constants are fixed before $h$ tends
to zero.  Equations \eqref{sgt:bose:number-law} and
\eqref{sgt:bose:finite-retention} then imply, for all sufficiently
small $h$,
\begin{equation}
 \Tr(Q\Xi_hN_{1,h}\Xi_h)\ge\frac{\mathsf n_Q}{2h},
 \qquad\Tr(Q\Xi_hN_{0,h}\Xi_h)=0.
 \label{sgt:bose:retained-gap}
\end{equation}
This proves a finite common receiver with bounded covariance and a
strictly extensive retained signal.

Suppose its actual effect has exact rejected-space exclusion and fixed
small error inside the retained space:
\begin{equation}
\begin{gathered}
 0\le E_h\le\Xi_h,\\
 \|E_h-\Xi_hQ\|\le e,
 \\
 e(\mathsf n_1+\mathsf n_0)<\mathsf n_Q/4.
 \label{sgt:bose:supported-effect}
\end{gathered}
\end{equation}
The inequality $0\le E_h\le\Xi_h$ implies
$E_h=\Xi_hE_h\Xi_h$ and
$\sqrt{E_h}=\Xi_h\sqrt{E_h}$.  Thus
\begin{equation}
 C_{i,h}=\eta\sqrt{E_h}N_{i,h}\sqrt{E_h}
   \le\eta n_bE_h,\qquad
 \|C_{i,h}\|\le\eta n_b.
 \label{sgt:bose:supported-covariance}
\end{equation}
Trace duality, \eqref{sgt:bose:retained-gap} and
\eqref{sgt:bose:number-law} leave a positive gap of order $h^{-1}$
between the two means.  Both means are $O(h^{-1})$.  The midpoint
threshold and \eqref{sgt:bose:concentration} now give
\begin{equation}
 P_{\mathrm{err}}^{\mathrm{filtered}}\le\ee^{-c/h}
 \label{sgt:bose:filtered-rate}
\end{equation}
for some $c>0$ and all sufficiently small $h$, also with an independent
calibrated Poisson dark mean at most $C_\nu/h$.  The exact support
condition in \eqref{sgt:bose:supported-effect} is substantive:
a small unrestricted effect error alone could admit bright rejected
modes and restore the larger covariance norm.  The following bound
quantifies a sufficient finite rejection tolerance.

\subsection{Finite amplitude leakage and the rate transition}
\label{sgt:bose:leakage}

Let the common passive receiver be a contraction
$T_h:\mathfrak h\to\mathfrak g_h$, where the counted output dimension
satisfies $\dim\mathfrak g_h\le D_*/h$ for a fixed $D_*$.  This explicit
output restriction keeps the finite-channel claim valid even when
infinitely many rejected input modes can leak.  Separate a fixed
$0<\eta\le1$ efficiency as before, and assume
\begin{equation}
\begin{gathered}
 E_{R,h}=\Xi_hT_h^*T_h\Xi_h,\\
 \|E_{R,h}-\Xi_hQ\|\le e,
 \\
 \ell_h=\|T_h(I-\Xi_h)\|.
 \label{sgt:bose:leakage-hypothesis}
\end{gathered}
\end{equation}
Set $A=T_h\Xi_h$ and $B=T_h(I-\Xi_h)$.  Since $\|A\|\le1$ and
$\|B\|=\ell_h$, the identity
$T_h^*T_h-E_{R,h}=A^*B+B^*A+B^*B$ gives
\begin{align}
 \|T_h^*T_h-\Xi_hQ\|&\le e+2\ell_h+\ell_h^2,\nonumber\\
 |\Tr[(T_h^*T_h-E_{R,h})N_{i,h}]|
    &\le(2\ell_h+\ell_h^2)\Tr N_{i,h}.
 \label{sgt:bose:leakage-mean}
\end{align}
For covariance, the positive-operator inequality
$(A+B)N(A+B)^*\le2ANA^*+2BNB^*$ follows by expanding
$(A-B)N(A-B)^*\ge0$.  The first term is bounded in norm by $n_b$
and the second by $\ell_h^2\|N_{i,h}\|$.  Therefore
\begin{equation}
 \|T_hN_{i,h}T_h^*\|
    \le2n_b+2\ell_h^2\|N_{i,h}\|
    \le2n_b+C\ell_h^2h^{-\gamma}.
 \label{sgt:bose:leakage-covariance}
\end{equation}
The covariance penalty is quadratic in rejected amplitude, although
the worst-case mean change includes linear cross terms.

For a precise sufficient uniform gap assumption, require
\begin{equation}
 \limsup_{h\downarrow0}(e+2\ell_h+\ell_h^2)
       <\frac{\mathsf n_Q}{4(\mathsf n_1+\mathsf n_0)}.
 \label{sgt:bose:leakage-gap}
\end{equation}
Together with \eqref{sgt:bose:retained-gap} and the number laws, this
implies a fixed positive $h^{-1}$ gap for the means of
$C_{i,h}=\eta T_hN_{i,h}T_h^*$, while both means remain $O(h^{-1})$.
The same is true after a common independent calibrated Poisson dark
count of mean at most $C_\nu/h$.  Using
\eqref{sgt:bose:leakage-covariance} in the common concentration bound
therefore proves, for some fixed $c>0$ and all sufficiently small $h$,
\begin{equation}
 P_{\mathrm{err}}^{\mathrm{leaky}}
  \le\exp\!\left[-\frac{c}
           {h(1+\ell_h^2h^{-\gamma})}\right].
 \label{sgt:bose:leakage-rate}
\end{equation}
Constants depend on the fixed profile, geometry, $\gamma,\tau,\eta$,
the chosen $b,R$, the gap margin and the dark-count bound.
If $e<\mathsf n_Q/[4(\mathsf n_1+\mathsf n_0)]$ and
$\ell_h=O(h^{\gamma/2})$, the gap condition is eventually satisfied
and the covariance is uniformly bounded, preserving
$\exp(-c/h)$.  For $\gamma=1/2$ this is the sufficient amplitude
tolerance $O(h^{1/4})$, equivalently a power leakage $O(h^{1/2})$.
Fixed sufficiently small leakage gives the available
$\exp(-c h^{\gamma-1})$ guarantee; intermediate leakage scales give
the displayed interpolation.  This is a sufficient receiver budget,
not a necessity theorem for every conceivable detector.

The common space $W_h$ uses both known candidate operators and never
depends on which hypothesis is true.  Its finite Fourier intersection
requires the corresponding spectral data and controlled projection
onto $T_{R,h}W_h$.  Rank $O(h^{-1})$ does not bound numerical conditioning,
settings precision, routing depth, stopband performance, physical mode
conversion, transmission, reset time or reservoir count.  Moreover a
finite detector marginal has covariance $\Xi_hN_{i,h}\Xi_h$.
It is not generally the Bose function of a compressed Hamiltonian:
compression and the nonlinear function $x\mapsto(\ee^{\beta x}-1)^{-1}$
do not commute.  The proof uses covariance compression throughout.

\subsection{Common vacuum, mechanical quantization and operational scope}
\label{sgt:bose:physical-scope}

The preceding laws are exact statements in the declared common
number-conserving model.  For finitely many real canonical coordinates
with $[x_j,p_k]=\mathrm i\delta_{jk}$ and a real symmetric positive
matrix $H$, the common annihilators $a=(x+\mathrm ip)/\sqrt2$ satisfy
\begin{equation}
 a^*Ha=\tfrac12(x^THx+p^THp-\Tr H).
 \label{sgt:bose:number-conserving-realization}
\end{equation}
Expansion proves the identity: the antisymmetric commutator contributes
$-\Tr H$ and the symmetric cross terms cancel.  It realizes $H$ in
both quadratures and conserves the chosen particle number.  It is a
model specification, not a claim that a geometric boundary operator
has available local couplings.

An ordinary mass-normalized mechanical Hamiltonian instead has the
form $(p^Tp+x^T\Omega^2x)/2$; its natural annihilators depend on
$\Omega$.  Already for one oscillator, let
\begin{equation}
 a_\omega=\sqrt{\omega/2}\,x+
                 \frac{\mathrm ip}{\sqrt{2\omega}},\qquad
 a_{\mathrm r}=\sqrt{\omega_{\mathrm r}/2}\,x+
                 \frac{\mathrm ip}{\sqrt{2\omega_{\mathrm r}}}.
 \label{sgt:bose:oscillator-modes}
\end{equation}
Solving for $x,p$ and substituting gives
\begin{equation}
 a_\omega=\cosh r\,a_{\mathrm r}
                 +\sinh r\,a_{\mathrm r}^*,
 \qquad r=\tfrac12\log(\omega/\omega_{\mathrm r}).
 \label{sgt:bose:bogoliubov}
\end{equation}
Inverting this relation shows that the $\omega$-vacuum has
$\langle a_{\mathrm r}^*a_{\mathrm r}\rangle=\sinh^2r$ quanta in the
fixed reference basis.  Its reference-vacuum probability is
\begin{equation}
 |\langle0_{\mathrm r}|0_\omega\rangle|^2
   =\frac{2\sqrt{\omega\omega_{\mathrm r}}}{\omega+\omega_{\mathrm r}}
   =\frac1{\cosh r}.
 \label{sgt:bose:vacuum-overlap}
\end{equation}
For a direct proof, integrate the normalized ground-state wavefunctions
$(\omega/\pi)^{1/4}\ee^{-\omega x^2/2}$ and
$(\omega_{\mathrm r}/\pi)^{1/4}\ee^{-\omega_{\mathrm r}x^2/2}$;
squaring their Gaussian overlap gives the displayed fraction.
Sending $\omega\to\infty$ suppresses thermal excitation in its own
natural basis but drives this fixed-basis vacuum probability to zero.
The conversion mixes creation and annihilation operators and is not
a passive rotation of a common annihilation basis.

Thus neither a common geometric $L^2$ space nor a growing mechanical
transverse stiffness establishes the reference-vacuum premise of
\eqref{sgt:bose:ideal-error} or \eqref{sgt:bose:reference-gap}.
A proposed platform must specify its canonical variables and common
detector modes and either realize the number-conserving model or
account for the creation part, noise, work and bandwidth of a
Bogoliubov conversion.  No infinite-mode unitary equivalence of Fock
representations is inferred here.  Ordinary displacement sensing,
passive mixing of layer coordinates and counting excitations in each
candidate's own normal modes are not automatically the same common
$Q$ occupation measurement.

Finally \eqref{sgt:bose:gaussian-mixture} supplies a mathematical
programmed realization of any finite output marginal: sample its
known classical complex Gaussian amplitudes and prepare the associated
coherent states.  This construction knows the covariance in advance.
It is a faithful encoded test of the stated distributions, not by
itself an unknown-specimen sensor or a computational quantum
advantage.  The geometric theorem, the controlled one-excitation
route and the full-Bose receiver establish distinct mathematical
access statements.  Their native physical implementation still
requires a geometry-to-generator approximation over the growing
band, a common mode map, and preparation, conversion and readout
satisfying the stated resource and noise budgets simultaneously.

\section{An exact finite circular counting certificate}
\label{sgt:finite:appendix}

\subsection{Finite statement and the declared ensemble}
\label{sgt:finite:certificate}

In the unit Dirichlet disk, compare two circular pushforward masses
of radii $r_1=99/200$ and $r_2=101/200$, each of mass $1/2$, with
the mass-one reference circle $r_0=1/2$. All are uniform in angle.
Use the common number-conserving Bose model with
$\gamma=1/2$, $\beta=1/20$; equivalently
$\varepsilon=1/100$, $\tau=1/2$.
The reference occupies $p=(1,1)^T/\sqrt2$ and its orthogonal layer
mode $q=(1,-1)^T/\sqrt2$ is vacuum. In every angular sector
$50\le|n|\le500$, detect the mode
$q_{\theta_n}=\cos\theta_n\,q+\sin\theta_n\,p$, where
$|\theta_n|\le1/10$ radians. The same settings apply under both
hypotheses, but may vary separately across all sectors. Detection
efficiency is $\eta=7/10$, and independent Poisson dark counts have
total gate mean $d\in[0,1/10]$. Declare the split state exactly when
the total recorded occupation $K$ is at least five. Then uniformly
over these allowances,
\begin{equation}
\begin{gathered}
 \Pr_0(K\ge5)<0.0057,\\
 \Pr_1(K\le4)<0.00072,\\
 P_{\rm err}=\tfrac12[\Pr_0(K\ge5)+\Pr_1(K\le4)]<0.0032.
 \label{sgt:finite:error-statement}
\end{gathered}
\end{equation}
There are $902$ measured output modes and $1804$ modes in the full
two-layer prepared band. This is a statement about the explicitly
aligned Bose model or its exact band marginal. It is not a
single-phonon error statement or a claim about an unmodified membrane.

\subsection{Exact Green blocks with both angular signs}
\label{sgt:finite:blocks}

For $z=re^{i\theta}$ and $z'=se^{i\phi}$ in the disk,
\[
 G(z,z')=\frac1{2\pi}\log
           \left|\frac{1-z\overline{z'}}{z-z'}\right|.
\]
Expanding $\log|1-te^{i\psi}|=-\sum_{n\ge1}t^n\cos(n\psi)/n$
for $|t|<1$ gives
\[
\begin{gathered}
 G=\frac1{2\pi}\bigl[-\log\max(r,s)
\\
 +\sum_{n\ge1}\frac{(\min(r,s)/\max(r,s))^n-(rs)^n}{n}
\\
                      \cos(n(\theta-\phi))\bigr].
\end{gathered}
\]
For $r=s$ the coefficients follow by the integrable logarithmic
limit. Convolution uses the normalized physical angular measure
$\dd\phi/(2\pi)$; a nonzero complex Fourier coefficient of the
cosine is $1/2$. After conjugating the masses, the $n$th block is
\begin{equation}
\begin{gathered}
 (K_n)_{ij}=\frac{\sqrt{\alpha_i\alpha_j}}{4\pi n}
 \left[\left(\frac{\min(r_i,r_j)}{\max(r_i,r_j)}\right)^n
                         -(r_ir_j)^n\right],\\
 n\ge1.
 \label{sgt:finite:green-block}
\end{gathered}
\end{equation}
Thus its prefactor here is $1/(8\pi n)$. The block for $-n$ is
identical. The reference nonzero scalar eigenvalue is
\begin{equation}
 k_{0,n}=\frac{1-(1/2)^{2n}}{4\pi n},\qquad
 K_{0,n}=k_{0,n}P.
 \label{sgt:finite:reference-block}
\end{equation}
Both total physical masses are one. Positivity of the Green form on
distinct layers gives positive blocks; the numerical certificate
additionally verifies their smaller eigenvalues are strictly positive
throughout this band. There are $451$ positive measured integers.
The independent $-n$ sectors, equivalently the second real harmonic
at each positive $n$, give multiplicity two. This is not a second
copy of the experiment. Rotational invariance makes full angular
blocks reducing for both candidates; their layer matrices need not
commute with each other.

\subsection{Covariance and sorter bounds}
\label{sgt:finite:sorter}

Let
\[
 g(x)=\frac1{e^{\beta/\sqrt x}-1},\qquad x>0,
 \quad N_n=g(K_n).
\]
Write $k_-<k_+$ for the eigenvalues of $K_n$,
$n_\pm=g(k_\pm)$, and $b_n=(n_+-n_-)/(k_+-k_-)>0$.
The exact two-dimensional functional calculus
\[
 g(K_n)=\frac{n_++n_-}{2}I+
 b_n\left(K_n-\frac{k_++k_-}{2}I\right)
\]
gives
\begin{equation}
\begin{gathered}
 q_\theta^*N_nq_\theta=
 \frac{n_++n_-}{2}
\\
 +b_n\bigl[\frac{K_{n,11}-K_{n,22}}2\sin(2\theta)
\\
                         -K_{n,12}\cos(2\theta)\bigr].
 \label{sgt:finite:exact-sorted-covariance}
\end{gathered}
\end{equation}
Since $K_{n,12}>0$ and $g$ is increasing, for $a=1/10$ the
detected target mean is bounded below, uniformly over $|\theta|\le a$,
by
\begin{equation}
\begin{gathered}
 \ell_n=\eta\bigl[\frac{n_++n_-}{2}-b_nK_{n,12}
\\
 -\frac{|b_n(K_{n,11}-K_{n,22})|}{2}\sin(2a)\bigr]>0.
 \label{sgt:finite:signal-lower}
\end{gathered}
\end{equation}
The inequality uses $\cos(2\theta)\le1$ in the negative term and
$|\sin(2\theta)|\le\sin(2a)$. These separate worst-case bounds need
not be simultaneously attained. Under the reference the detected
mean is exactly $\eta g(k_{0,n})\sin^2\theta$, and therefore at most
\begin{equation}
 u_n=\eta g(k_{0,n})\sin^2a.
 \label{sgt:finite:null-upper}
\end{equation}

\subsection{Geometric occupation tails and interval arithmetic}
\label{sgt:finite:tails}

A single mode in a gauge-invariant Gaussian state with mean $\nu$ has
geometric occupation distribution
$\Pr(X=k)=\nu^k/(1+\nu)^{k+1}$, so
\[
 \mathbb Ez^X=[1+\nu(1-z)]^{-1}.
\]
One direct verification uses its positive coherent-state mixture:
draw a circular complex Gaussian amplitude with variance $\nu$,
then average the coherent-state Poisson generating function
$\exp[-(1-z)|\zeta|^2]$. Loss of efficiency $\eta$ rescales the
variance by $\eta$. Distinct angular sectors are independent
conditional on the fixed sorter settings; no independence of the
two layer eigenmodes after sorting is assumed or needed. The same
count structure is classical thermal optics \cite{Lu2018,Klich2002}.

For null tilt $t=3/2$, each denominator
$1-u_n(e^t-1)$ is positive. Markov's inequality, monotonicity in the
geometric mean and the Poisson generating function imply
\begin{equation}
\begin{gathered}
 B_0=\exp\bigl[-5t-2\sum_{n=50}^{500}
\\
             \log(1-u_n(e^t-1))\\
+\tfrac1{10}(e^t-1)\bigr]
\\
 \ \ge\Pr_0(K\ge5).
 \label{sgt:finite:null-chernoff}
\end{gathered}
\end{equation}
For the miss event $K\le4$, choose $s=1$. Dropping nonnegative
independent dark counts increases its probability. Its decreasing
Laplace transform and \eqref{sgt:finite:signal-lower} give
\begin{equation}
\begin{gathered}
 B_1=\exp\bigl[4s-2\sum_{n=50}^{500}
\\
              \log(1+\ell_n(1-e^{-s}))\bigr]
\\
 \ \ge\Pr_1(K\le4).
 \label{sgt:finite:signal-chernoff}
\end{gathered}
\end{equation}
The coefficient $4$, rather than $5$, is the largest integer in the
miss event. These bounds hold conditional on every allowed collection
of angles, so they also hold after any classical randomization of
the angles independent of the prepared state, even if the angles
are correlated across sectors.

The accompanying checker evaluates these finite expressions with
Arb ball arithmetic \cite{Johansson2017}. Every input rational is
formed by exact integer division; $\pi$, roots, exponentials, sines
and logarithms receive interval enclosures. For every $n$ it checks
$k_->0$, $k_+-k_->0$, $b_n>0$, $K_{n,12}>0$, $\ell_n>0$, and
$1-u_n(e^t-1)>0$. Explicitly, the discriminant is computed as
\[
 \sqrt{(K_{n,11}-K_{n,22})(K_{n,11}-K_{n,22})+4K_{n,12}K_{n,12}},
\]
so a tiny zero-straddling diagonal difference is handled by valid
interval multiplication. Finally it proves, as strict interval
inequalities against exact rationals,
\[
 B_0<\frac{57}{10000},\qquad
 B_1<\frac{72}{100000},\qquad
 \frac{B_0+B_1}{2}<\frac{32}{10000}.
\]
At 256-bit working precision the enclosures lie inside
\begin{gather*}
 B_0\in(0.00566995032088188856,\\
 \hspace{2em}0.00566995032088188858),\\
 B_1\in(0.00071698918995924630,\\
 \hspace{2em}0.00071698918995924633),\\
 (B_0+B_1)/2\in(0.00319346975542056743,\\
 \hspace{2em}0.00319346975542056746).
\end{gather*}
These finite rational comparisons prove
\eqref{sgt:finite:error-statement}. Higher-precision replay checks
the same expressions. The quantities $B_0,B_1$ are upper bounds on
tails, not exact tail probabilities or Monte Carlo confidence limits.

\subsection{Resource values and exact scope of the marginal}
\label{sgt:finite:scope}

The full prepared band occupation and excitation-energy sums are
\begin{gather*}
 \overline N_1=2\sum_{n=50}^{500}(n_-+n_+),\\
 \overline E_1=2\sum_{n=50}^{500}
                  (n_-/\sqrt{k_-}+n_+/\sqrt{k_+}),\\
 \overline N_0=2\sum_{n=50}^{500}g(k_{0,n}),\\
 \overline E_0=2\sum_{n=50}^{500}g(k_{0,n})/\sqrt{k_{0,n}}.
\end{gather*}
Their interval-enclosed values, rounded here for interpretation, are
\[
\begin{gathered}
 (\overline N_1,\overline N_0)\simeq(61.79306,81.31572),\\
 (\overline E_1,\overline E_0)\simeq(3550.20358,3590.80726).
\end{gathered}
\]
These are pre-detection band resources and omit zero-point energy.
In contrast, the uniform detected target mean lower bound is about
$18.00837$, and the reference upper bound before dark counts is about
$0.56732$. A complete equilibrium sample has additional occupation
and energy outside the band. The unequal prepared occupations are
not equal thermal scalar data; this finite point certifies access
and discrimination, not the asymptotic thermal-scalar comparison or
an advantage over every receiver at matched brightness.

The full two-layer angular band is an exact finite reducing subspace.
Within each block the detected rank-one output need not reduce the
Hamiltonian. Its marginal is nevertheless Gaussian with the exact
compressed covariance used in
\eqref{sgt:finite:exact-sorted-covariance}. Tracing out the remaining
modes does not assert that they are unoccupied and introduces no
spectral-tail approximation into this measurement. One can prepare
the specified band marginal alone, for example by sampling correlated
classical complex Gaussian amplitudes with its covariance. This is
a programmed simulation with known covariance. In general the
covariance $V^*NV$ must not be replaced by
$(e^{\beta V^*HV}-I)^{-1}$, because compression and functional
calculus do not commute. Reproducing this marginal is not global
Fock-state approximation in trace distance.

The finite certificate assumes ideal exclusion of angular modes
outside $50\le|n|\le500$, the stated efficiency throughout the band,
real within-sector rotation errors only, number resolution through
the threshold, and the independent gate-level Poisson background.
It gives no finite numerical allowance for stopband leakage,
cross-sector mixing, phase errors, finite radial mass width, changing
carrier stiffness, or bulk inertia. The general asymptotic rejection
theorem does not by itself supply those finite tolerances.
The common vacuum and annihilation modes are declared preparation
assumptions. The classical mode-projection receiver and its
crosstalk sensitivity have direct precedents
\cite{Lu2018,Schlichtholz2024}; the result here is this exact
geometric finite instance with its resource and tolerance ledger.

% Primary sources and explicitly identified books/surveys.
\begin{thebibliography}{99}

\bibitem{KondejKrejcirik2013}
S. Kondej and D. Krej\v{c}i\v{r}\'{\i}k,
Spectral analysis of a quantum system with a double line singular interaction,
\href{https://doi.org/10.4171/PRIMS/121}
{Publ. Res. Inst. Math. Sci. \textbf{49}, 831--859 (2013)}.

\bibitem{KondejKrejcirik2016}
S. Kondej and D. Krej\v{c}i\v{r}\'{\i}k,
Asymptotic spectral analysis in colliding leaky quantum layers,
\href{https://arxiv.org/abs/1606.07273}
{J. Math. Anal. Appl. \textbf{446}, 1328--1355 (2017)}
(preprint 2016).

\bibitem{HundertmarkLiebThomas1998}
D. Hundertmark, E. H. Lieb, and L. E. Thomas,
A sharp bound for an eigenvalue moment of the one-dimensional
Schr\"odinger operator,
\href{https://arxiv.org/abs/math-ph/9806012}
{Adv. Theor. Math. Phys. \textbf{2}, 719--731 (1998)}.

\bibitem{DallaRivaProvenzano2018}
M. Dalla Riva and L. Provenzano,
On vibrating thin membranes with mass concentrated near the boundary:
An asymptotic analysis,
\href{https://doi.org/10.1137/17M1118221}
{SIAM J. Math. Anal. \textbf{50}, 2928--2967 (2018)}.

\bibitem{Golovaty2022}
Y. Golovaty,
Membranes with thin and heavy inclusions: Asymptotics of spectra,
\href{https://doi.org/10.3233/ASY-211743}
{Asymptot. Anal. \textbf{130}, 23--51 (2022)}
(published online 2021; \href{https://arxiv.org/abs/2101.10059}{arXiv:2101.10059}).

\bibitem{GolovatyGomezPerezMartinez2025}
Y. Golovaty, D. G\'omez, and M.-E. P\'erez-Mart\'inez,
On eigenvibrations of branched structures with heterogeneous mass density,
\href{https://doi.org/10.1016/j.jmaa.2025.129586}
{J. Math. Anal. Appl. \textbf{549}, 129586 (2025)}.

\bibitem{AssalDimassiFujii2019}
M. Assal, M. Dimassi, and S. Fujii\'{e},
Semiclassical trace formula and spectral shift function for systems
via a stationary approach,
\href{https://doi.org/10.1093/imrn/rnx149}
{Int. Math. Res. Not. \textbf{2019}, 1227--1264 (2019)}.

\bibitem{Kuster2017}
B. K\"uster,
On the semiclassical functional calculus for $h$-dependent functions,
\href{https://doi.org/10.1007/s10455-017-9549-1}
{Ann. Global Anal. Geom. \textbf{52}, 57--97 (2017)}.

\bibitem{BollmannMuller2024}
L. Bollmann and P. M\"uller,
The Widom--Sobolev formula for discontinuous matrix-valued symbols,
\href{https://doi.org/10.1016/j.jfa.2024.110651}
{J. Funct. Anal. \textbf{287}, 110651 (2024)}.

\bibitem{Lu2018}
X.-M. Lu, H. Krovi, R. Nair, S. Guha, and J. H. Shapiro,
Quantum-optimal detection of one-versus-two incoherent optical sources
with arbitrary separation,
\href{https://doi.org/10.1038/s41534-018-0114-y}
{npj Quantum Inf. \textbf{4}, 64 (2018)}.

\bibitem{Schlichtholz2024}
K. Schlichtholz, T. Linowski, M. Walschaers, N. Treps,
\L{}. Rudnicki, and G. Sorelli,
Practical tests for sub-Rayleigh source discriminations with imperfect
demultiplexers,
\href{https://doi.org/10.1364/OPTICAQ.502459}
{Optica Quantum \textbf{2}, 29--34 (2024)}.

\bibitem{Liu2026}
J.-Q. Liu, C.-N. Hu, J. Xin, and X.-M. Lu,
Noise-robust discrimination of incoherent point sources with
spatial-mode demultiplexing,
\href{https://arxiv.org/abs/2608.01765}
{arXiv:2608.01765 (2026)}.

\bibitem{Klich2002}
I. Klich,
Full counting statistics: An elementary derivation of Levitov's formula,
\href{https://arxiv.org/abs/cond-mat/0209642}{arXiv:cond-mat/0209642 (2002)}.

\bibitem{Clements2016}
W. R. Clements, P. C. Humphreys, B. J. Metcalf,
W. S. Kolthammer, and I. A. Walmsley,
Optimal design for universal multiport interferometers,
\href{https://doi.org/10.1364/OPTICA.3.001460}
{Optica \textbf{3}, 1460--1465 (2016)}.

\bibitem{Lostaglio2015}
M. Lostaglio, D. Jennings, and T. Rudolph,
Description of quantum coherence in thermodynamic processes requires
constraints beyond free energy,
\href{https://doi.org/10.1038/ncomms7383}
{Nat. Commun. \textbf{6}, 6383 (2015)}.

\bibitem{Watrous2018}
J. Watrous,
\emph{The Theory of Quantum Information}
(Cambridge University Press, Cambridge, 2018), Theorem 3.4;
\href{https://cs.uwaterloo.ca/~watrous/TQI/TQI.pdf}{author's text}.

\bibitem{Simon1976}
B. Simon,
The bound state of weakly coupled Schr\"odinger operators in one and
two dimensions,
\href{https://doi.org/10.1016/0003-4916(76)90038-5}
{Ann. Phys. \textbf{97}, 279--288 (1976)}.

\bibitem{Simon2022}
B. Simon,
Twelve tales in mathematical physics: An expanded Heineman prize lecture,
\href{https://doi.org/10.1063/5.0056008}
{J. Math. Phys. \textbf{63}, 021101 (2022)}.

\bibitem{Quesada2018}
N. Quesada, J. M. Arrazola, and N. Killoran,
Gaussian boson sampling using threshold detectors,
\href{https://doi.org/10.1103/PhysRevA.98.062322}
{Phys. Rev. A \textbf{98}, 062322 (2018)}.

\bibitem{BruneauDerezinski2007}
L. Bruneau and J. Derezi\'{n}ski,
Bogoliubov Hamiltonians and one-parameter groups of Bogoliubov transformations,
\href{https://bruneau.perso.math.cnrs.fr/bogoliubovJMP.pdf}
{J. Math. Phys. \textbf{48}, 022101 (2007)}.

\bibitem{Aspelmeyer2014}
M. Aspelmeyer, T. J. Kippenberg, and F. Marquardt,
Cavity optomechanics,
\href{https://doi.org/10.1103/RevModPhys.86.1391}
{Rev. Mod. Phys. \textbf{86}, 1391--1452 (2014)}.

\bibitem{Garrigos2015}
G. Garrig\'{o}s,
A weak 2-weight problem for the Poisson--Hermite semigroup,
in \emph{Advanced Courses of Mathematical Analysis VI}
(World Scientific, 2017), pp. 153--171,
\href{https://doi.org/10.1142/9789813147645_0006}
{doi:10.1142/9789813147645\_0006};
\href{https://webs.um.es/gustavo.garrigos/papers/garrigos_antequera_proc_june2015.pdf}
{author's 2015 draft}, Eq. (1.1).

\bibitem{PlenioKnight1998}
M. B. Plenio and P. L. Knight,
The quantum-jump approach to dissipative dynamics in quantum optics,
\href{https://doi.org/10.1103/RevModPhys.70.101}
{Rev. Mod. Phys. \textbf{70}, 101--144 (1998)}.

\bibitem{DorflerBullo2013}
F. D\"orfler and F. Bullo,
Kron reduction of graphs with applications to electrical networks,
\href{https://doi.org/10.1109/TCSI.2012.2215780}
{IEEE Trans. Circuits Syst. I \textbf{60}, 150--163 (2013)}.

\bibitem{Johansson2017}
F. Johansson,
Arb: Efficient arbitrary-precision midpoint-radius interval arithmetic,
\href{https://doi.org/10.1109/TC.2017.2690633}
{IEEE Trans. Comput. \textbf{66}, 1281--1292 (2017)}.

\end{thebibliography}

\end{document}
